\documentclass[11pt,a4paper]{report}
\usepackage[T1]{fontenc}
\usepackage[utf8]{inputenc}
\usepackage[french]{babel}
\usepackage{amsmath,amssymb}
\usepackage{mathpazo}
\usepackage[scaled=0.92]{helvet}
\usepackage{icomma}
\usepackage[version=4]{mhchem}
\linespread{1.06}

% ------------------------------------------------------------------ Mise en page
\usepackage[a4paper,left=2.4cm,right=2.4cm,top=2.7cm,bottom=2.6cm,headsep=0.75cm,footskip=1.1cm]{geometry}
\setlength{\parindent}{1.5em}
\setlength{\parskip}{0.35em plus 0.15em minus 0.1em}

% ------------------------------------------------------------------ Couleurs
\usepackage{xcolor}
\definecolor{navy}{HTML}{193C78}
\definecolor{cdef}{HTML}{193C78}   \definecolor{cdefbg}{HTML}{F7F7FF}
\definecolor{cres}{HTML}{961E1E}   \definecolor{cresbg}{HTML}{FFF7F7}
\definecolor{cmet}{HTML}{1E6E32}   \definecolor{cmetbg}{HTML}{F7FFF7}
\definecolor{cex}{HTML}{4C4C4C}    \definecolor{cexbg}{HTML}{FAFAFA}
\definecolor{cog}{HTML}{642882}    \definecolor{cogbg}{HTML}{FAF8FC}
\definecolor{cpi}{HTML}{C86400}    \definecolor{cpibg}{HTML}{FCF7F2}
\definecolor{chp}{HTML}{0F6E78}    \definecolor{chpbg}{HTML}{F3FAFA}
\definecolor{crap}{HTML}{4F6688}   \definecolor{crapbg}{HTML}{F6F8FB}
\definecolor{c1}{HTML}{193C78}
\definecolor{c2}{HTML}{C86400}
\definecolor{c3}{HTML}{961E1E}
\definecolor{c4}{HTML}{1E6E32}
\definecolor{c5}{HTML}{642882}

% ------------------------------------------------------------------ Outils
\usepackage{etoolbox}
\usepackage[locale=FR,inter-unit-product=\ensuremath{{}\cdot{}},per-mode=power]{siunitx}
\usepackage{booktabs,array,tabularx,multirow}
\usepackage{enumitem}
\setlist[itemize]{label=\textemdash,leftmargin=2.2em,itemsep=0.2em,topsep=0.3em}
\setlist[enumerate]{leftmargin=2.2em,itemsep=0.2em,topsep=0.3em}
\usepackage[font=small,labelfont={sf,bf,color=navy},labelsep=endash,skip=6pt]{caption}
\usepackage{float}
\usepackage{needspace}
\usepackage{xltabular}
\usepackage{listings}
\makeatletter
\setlength{\@fptop}{0pt}
\setlength{\@fpsep}{2.2em}
\setlength{\@fpbot}{0pt plus 1fil}
\makeatother
\renewcommand{\floatpagefraction}{0.6}
\renewcommand{\topfraction}{0.85}
\renewcommand{\textfraction}{0.12}

\usepackage{tikz}
\usetikzlibrary{arrows.meta,positioning,calc,decorations.pathreplacing,patterns,shapes.geometric,babel}
\usepackage{pgfplots}
\pgfplotsset{compat=1.18}
\usepgfplotslibrary{fillbetween}
\pgfkeys{/pgf/number format/.cd,use comma,1000 sep={\,}}
\pgfplotsset{cours/.style={width=0.82\linewidth,height=0.46\linewidth,
  axis lines=left,tick label style={font=\small},label style={font=\small},
  every axis plot/.append style={line width=1.1pt},
  grid=major,grid style={black!12},legend style={font=\small,draw=black!25,fill=white}}}

% ------------------------------------------------------------------ Titres
\usepackage{titlesec}
\titleformat{\chapter}[display]
  {\sffamily\bfseries\color{navy}}
  {\normalsize\MakeUppercase{\chaptertitlename}~\thechapter}
  {0.35em}
  {\huge\raggedright}
  [\vspace{0.25em}{\color{navy}\titlerule[0.9pt]}]
\titleformat{name=\chapter,numberless}[block]
  {\sffamily\bfseries\color{navy}}{}{0pt}{\huge\raggedright}
  [\vspace{0.25em}{\color{navy}\titlerule[0.9pt]}]
\titlespacing*{\chapter}{0pt}{0.6cm}{1.3cm}
\titleformat{\section}{\sffamily\bfseries\Large\color{navy}}{\thesection}{0.7em}{}
\titlespacing*{\section}{0pt}{1.5em plus 0.4em minus 0.2em}{0.7em}
\titleformat{\subsection}{\sffamily\bfseries\large\color{navy}}{\thesubsection}{0.7em}{}
\titlespacing*{\subsection}{0pt}{1.1em plus 0.3em minus 0.2em}{0.4em}
\setcounter{secnumdepth}{2}
\setcounter{tocdepth}{1}

% ------------------------------------------------------------------ En-têtes
\usepackage{fancyhdr}
\pagestyle{fancy}
\fancyhf{}
\renewcommand{\chaptermark}[1]{\markboth{\thechapter.~#1}{}}
\fancyhead[L]{\sffamily\small\color{black!60}\nouppercase{\leftmark}}
\fancyhead[R]{\sffamily\small\color{black!60}Thermochimie~--~PT*}
\fancyfoot[C]{\sffamily\small\thepage}
\renewcommand{\headrulewidth}{0.5pt}
\setlength{\headheight}{14pt}
\fancypagestyle{plain}{\fancyhf{}\fancyfoot[C]{\sffamily\small\thepage}\renewcommand{\headrulewidth}{0pt}}

% ------------------------------------------------------------------ Encadrés
\usepackage[most]{tcolorbox}
\tcbset{cours/.style={enhanced,breakable,arc=1.2mm,boxrule=0.5pt,
  left=3mm,right=3mm,top=2mm,bottom=2mm,
  fonttitle=\sffamily\bfseries\small,coltitle=white,toptitle=0.7mm,bottomtitle=0.7mm,
  before skip=1em plus 0.2em,after skip=1em plus 0.2em,
  before upper={\setlength{\parindent}{0pt}\setlength{\parskip}{0.45em}}}}
\newcommand{\tsep}[2]{#1\ifstrempty{#2}{}{~\textemdash{}~#2}}
\newtcolorbox{definition}[1][]{cours,colframe=cdef,colback=cdefbg,title={\tsep{Définition}{#1}}}
\newtcolorbox{resultat}[1]{cours,colframe=cres,colback=cresbg,title={#1}}
\newtcolorbox{methode}[1][]{cours,colframe=cmet,colback=cmetbg,title={\tsep{Méthode}{#1}}}
\newtcolorbox{exemple}[1][]{cours,colframe=cex,colback=cexbg,title={\tsep{Exemple}{#1}}}
\newtcolorbox{ordres}[1][]{cours,colframe=cog,colback=cogbg,title={\tsep{Ordres de grandeur}{#1}}}
\newtcolorbox{pieges}{cours,colframe=cpi,colback=cpibg,title={Pièges classiques}}
\newtcolorbox{rappel}[1][]{cours,colframe=crap,colback=crapbg,title={\tsep{Rappel de première année}{#1}}}
\newtcolorbox{horsprog}[1][]{cours,colframe=chp,colback=chpbg,
  borderline={0.5pt}{0pt}{chp,dashed},
  title={\tsep{Hors programme}{#1}}}
\newtcolorbox{numerique}[1][]{cours,colframe=c5,colback=cogbg,title={\tsep{Capacité numérique}{#1}}}

\newenvironment{demo}[1][Démonstration]{\par\smallskip\noindent{\color{navy}\sffamily\bfseries\small #1.}~}{\hfill$\square$\par\smallskip}

% Exercices
\newcounter{exo}
\newcommand{\diff}[1]{\textcolor{c2}{\ifcase#1\or$\bigstar$\or$\bigstar\bigstar$\or$\bigstar\bigstar\bigstar$\fi}}
\newcommand{\exercice}[2]{\refstepcounter{exo}\needspace{6\baselineskip}\section*{Exercice~\theexo~--~#1\hfill\diff{#2}}\addcontentsline{toc}{section}{Exercice~\theexo~--~#1}}
\newcommand{\correction}{\needspace{4\baselineskip}\par\medskip\noindent{\color{c4}\sffamily\bfseries Correction}\par\nopagebreak}

% ------------------------------------------------------------------ Code Python
\lstset{language=Python,basicstyle=\ttfamily\footnotesize,keywordstyle=\color{navy}\bfseries,
  commentstyle=\color{black!55}\itshape,stringstyle=\color{c3},showstringspaces=false,
  columns=fullflexible,keepspaces=true,frame=none,
  literate={é}{{\'e}}1 {è}{{\`e}}1 {à}{{\`a}}1 {ê}{{\^e}}1 {É}{{\'E}}1 {ξ}{{$\xi$}}1 {°}{{\textdegree}}1}

% ------------------------------------------------------------------ Macros
\newcommand{\dd}{\mathrm{d}}
\newcommand{\rG}{\Delta_r G}
\newcommand{\rH}{\Delta_r H}
\newcommand{\rS}{\Delta_r S}
\newcommand{\rGo}{\Delta_r G^\circ}
\newcommand{\rHo}{\Delta_r H^\circ}
\newcommand{\rSo}{\Delta_r S^\circ}
\newcommand{\fHo}{\Delta_f H^\circ}
\newcommand{\fGo}{\Delta_f G^\circ}
\newcommand{\Ko}{K^\circ}
\newcommand{\Qr}{Q_r}
\newcommand{\Po}{P^\circ}
\newcommand{\co}{c^\circ}
\newcommand{\sumi}{\sum_i}
\newcommand{\kJmol}{\,\mathrm{kJ\cdot mol^{-1}}}
\newcommand{\JKmol}{\,\mathrm{J\cdot K^{-1}\cdot mol^{-1}}}


\begin{document}

% ================================================================== Page de titre
\begin{titlepage}
\centering
\vspace*{3.2cm}
{\sffamily\color{navy}\small\MakeUppercase{Physique-chimie -- PT*}\par}
\vspace{0.8cm}
{\color{navy}\rule{\linewidth}{0.9pt}\par}
\vspace{0.6cm}
{\sffamily\bfseries\color{navy}\Huge Thermodynamique de la\\[0.3em] transformation chimique\par}
\vspace{0.5cm}
{\large Grandeurs de réaction, enthalpie libre, équilibres,\\ thermodynamique de l'oxydoréduction et piles\par}
\vspace{0.6cm}
{\color{navy}\rule{\linewidth}{0.9pt}\par}
\vspace{1.6cm}
\begin{tikzpicture}
  \begin{axis}[width=8.5cm,height=5cm,axis lines=left,xtick=\empty,ytick=\empty,
    xlabel={$\xi$},ylabel={$G$},xmin=0,xmax=1.05,ymin=-3.4,ymax=0.7,clip=false,
    label style={font=\small},every axis plot/.append style={line width=1.3pt}]
    \addplot[c1] table[x=xi,y=G]{data/gxi.dat};
    \addplot[c3,dashed,line width=0.8pt] coordinates {(0.691,-3.4) (0.691,-2.915)};
    \addplot[only marks,mark=*,mark size=2pt,c3] coordinates {(0.691,-2.915)};
    \node[c3,font=\small,anchor=west] at (axis cs:0.71,-3.2) {équilibre : $\Delta_rG=0$};
  \end{axis}
\end{tikzpicture}
\vfill
{\small\color{black!65} Cours complet conforme au programme officiel de PT (2021)\\
parties « Transformations chimiques de la matière : aspects thermodynamiques »\\
et « Étude thermodynamique des réactions d'oxydoréduction -- piles »\par}
\vspace{0.5cm}
\end{titlepage}

\tableofcontents

% ================================================================== Chapitre 0
\chapter*{Programme, plan et notations}
\addcontentsline{toc}{chapter}{Programme, plan et notations}
\markboth{Programme, plan et notations}{}

\section*{Ce que le programme exige}

Ce chapitre regroupe trois blocs du programme de PT. Chaque capacité ci-dessous est traitée dans le cours (le numéro du chapitre est indiqué entre crochets) et au moins un exercice corrigé la mobilise.

\begin{resultat}{Capacités exigibles -- résumé}
\textbf{Premier principe appliqué à la transformation chimique}
\begin{itemize}
  \item calculer une enthalpie standard de réaction avec des tables ou la loi de Hess ; relier son signe au caractère endo- ou exothermique \hfill[3]
  \item prévoir le sens et estimer le transfert thermique en réacteur monobare isotherme \hfill[4]
  \item évaluer la température atteinte en réacteur adiabatique monobare ; simuler $T(t)$ en Python \hfill[4]
\end{itemize}
\textbf{Deuxième principe appliqué à la transformation chimique}
\begin{itemize}
  \item prévoir le signe de l'entropie standard de réaction \hfill[3]
  \item relier création d'entropie et enthalpie libre de réaction à $T$ et $P$ fixées \hfill[5]
  \item prévoir le sens d'évolution avec $\Delta_r G$ (ou $Q_r$ et $K^\circ$) \hfill[5]
  \item calculer les grandeurs standard de réaction (tables, Hess, combinaisons linéaires) \hfill[3]
  \item citer et exploiter la relation de Van~'t~Hoff ; calculer $K^\circ(T)$ à toute température \hfill[6]
  \item déterminer l'état final : équilibre chimique ou transformation totale \hfill[5]
  \item identifier les paramètres d'optimisation d'une synthèse \hfill[6]
\end{itemize}
\textbf{Thermodynamique de l'oxydoréduction et piles}
\begin{itemize}
  \item relier $\Delta_rG$ aux potentiels des couples et $\Delta_rG^\circ$ aux potentiels standard \hfill[7]
  \item calculer un potentiel standard inconnu à partir de données thermodynamiques \hfill[7]
  \item établir l'inégalité entre variation d'enthalpie libre et travail électrique \hfill[8]
  \item relier la tension à vide d'une pile à l'enthalpie libre de réaction ; calculer sa capacité \hfill[8]
\end{itemize}
\end{resultat}

\noindent Trois précisions du programme guident tout le cours :
\begin{itemize}
  \item l'expression $\mu_i=\mu_i^\circ+RT\ln a_i$ du potentiel chimique est \emph{admise} (sa démonstration est hors programme) ;
  \item on suppose $\rHo$ et $\rSo$ \emph{indépendantes de la température} en dehors des changements d'état (approximation d'Ellingham) ;
  \item la notion d'\emph{affinité chimique} n'est pas utilisée : le sens d'évolution se lit sur le signe de $\Delta_rG$.
\end{itemize}

\section*{Notations utilisées}

\renewcommand{\arraystretch}{1.25}
\begin{xltabular}{\linewidth}{>{\centering\arraybackslash}p{2.2cm} X >{\centering\arraybackslash}p{3.2cm}}
\toprule
\textbf{Symbole} & \textbf{Signification} & \textbf{Unité} \\ \midrule
$\nu_i$ & coefficient stœchiométrique \emph{algébrique} ($>0$ produit, $<0$ réactif) & sans \\
$\xi$ & avancement de la réaction & mol \\
$a_i$ & activité du constituant $\mathrm{A}_i$ & sans \\
$\mu_i$, $\mu_i^\circ$ & potentiel chimique, potentiel chimique standard & $\mathrm{J\cdot mol^{-1}}$ \\
$G$, $H$, $S$ & enthalpie libre, enthalpie, entropie du système & J, J, $\mathrm{J\cdot K^{-1}}$ \\
$\Delta_r X$ & grandeur de réaction ($X = G, H, S$) & $\mathrm{J\cdot mol^{-1}}$ (ou $\mathrm{J\cdot K^{-1}\cdot mol^{-1}}$) \\
$\Delta_r X^\circ$ & grandeur standard de réaction & idem \\
$\Delta_f X^\circ$ & grandeur standard de formation & idem \\
$Q_r$ & quotient de réaction & sans \\
$K^\circ$ & constante thermodynamique d'équilibre (notée aussi $K$) & sans \\
$Q$ & transfert thermique reçu par le système & J \\
$\gamma$ & nombre d'électrons échangés dans la réaction redox & sans \\
$E$, $E^\circ$ & potentiel d'électrode (de Nernst), potentiel standard & V \\
$e$, $e^\circ$ & tension à vide d'une pile ($e=E_+-E_-$), tension à vide standard & V \\
$q$, $q_{\max}$ & charge débitée, capacité d'une pile & C (ou A$\cdot$h) \\
\bottomrule
\end{xltabular}

\begin{ordres}[constantes utiles]
\begin{itemize}
  \item constante des gaz parfaits : $R = \SI{8,314}{J.K^{-1}.mol^{-1}}$ ; constante de Faraday : $F = \mathcal{N}_A e = \SI{96485}{C.mol^{-1}} \approx \SI{96500}{C.mol^{-1}}$ ;
  \item pression standard : $P^\circ = \SI{1}{bar} = \SI{e5}{Pa}$ ; concentration standard : $c^\circ = \SI{1}{mol.L^{-1}}$ ;
  \item à $T = \SI{298}{K}$ : $RT = \SI{2,48}{kJ.mol^{-1}}$ et
  \[ \frac{RT\ln 10}{F} = \SI{0,0592}{V} \approx \SI{0,06}{V} \qquad\text{(c'est le « 0,06 » des formules de Nernst).} \]
\end{itemize}
\end{ordres}

% ================================================================== Chapitre 1
\chapter{Rappels de première année}

Ce chapitre rassemble, sous forme condensée, les outils de PTSI réutilisés en permanence ensuite : avancement, activité, quotient de réaction et oxydoréduction. Les diagrammes $E$-pH et les titrages rédox, également revus en début de chapitre, ne sont pas repris ici : ils n'interviennent pas dans les raisonnements thermodynamiques qui suivent.

\section{Système physico-chimique et avancement}

\begin{definition}[système physico-chimique, constituant, phase]
\begin{itemize}
  \item Un \textbf{constituant} est une espèce chimique \emph{dans un état physique donné} : \ce{H2O(l)} et \ce{H2O(g)} sont deux constituants différents.
  \item Une \textbf{phase} est une région de l'espace où les grandeurs intensives (masse volumique, concentrations\dots) varient continûment. Un mélange de gaz forme toujours une seule phase ; deux solides distincts forment deux phases.
  \item Un \textbf{système physico-chimique} est un ensemble de constituants répartis en une ou plusieurs phases, susceptible d'être le siège de transformations chimiques ou physiques.
\end{itemize}
\end{definition}

\begin{definition}[équation de réaction, coefficients stœchiométriques algébriques]
On écrit toute équation de réaction sous la forme
\[ 0 = \sum_i \nu_i\,\mathrm{A}_i \]
où $\nu_i$ est le \textbf{coefficient stœchiométrique algébrique} du constituant $\mathrm{A}_i$ : $\nu_i>0$ pour un produit, $\nu_i<0$ pour un réactif, $\nu_i = 0$ pour une espèce spectatrice.
\end{definition}

\begin{exemple}
Pour \ce{N2(g) + 3 H2(g) = 2 NH3(g)} : $\nu_{\ce{N2}} = -1$, $\nu_{\ce{H2}} = -3$, $\nu_{\ce{NH3}} = +2$.
\end{exemple}

\begin{definition}[avancement]
L'\textbf{avancement} $\xi$ (en mol) est défini par
\[ n_i(\xi) = n_{i,0} + \nu_i\,\xi \qquad\Longleftrightarrow\qquad \dd n_i = \nu_i\,\dd\xi \quad\text{pour tout }i. \]
Il est positif si la réaction avance dans le sens direct ($\rightarrow$), négatif dans le sens indirect ($\leftarrow$). L'\textbf{avancement maximal} $\xi_{\max}$ correspond à l'épuisement du réactif limitant.
\end{definition}

\section{Activité, quotient de réaction et constante d'équilibre}

\begin{definition}[activité]
L'\textbf{activité} $a_i$ d'un constituant est une grandeur sans dimension qui mesure son « poids » effectif dans le mélange. Dans le cadre du programme (cas idéaux) :

\renewcommand{\arraystretch}{1.3}
\begin{center}
\begin{tabular}{ll}
\toprule
\textbf{Constituant} & \textbf{Activité} \\ \midrule
gaz parfait dans un mélange & $a_i = \dfrac{P_i}{P^\circ} = x_i\dfrac{P}{P^\circ}$ \quad ($P_i$ pression partielle) \\[0.6em]
soluté en solution diluée & $a_i = \dfrac{c_i}{c^\circ}$ \\[0.6em]
solvant (eau) & $a = 1$ \\
solide ou liquide \emph{pur} (seul dans sa phase) & $a = 1$ \\
\bottomrule
\end{tabular}
\end{center}
\end{definition}

\begin{definition}[quotient de réaction]
Pour la réaction $0 = \sum_i \nu_i\mathrm{A}_i$, le \textbf{quotient de réaction} est
\[ Q_r = \prod_i a_i^{\,\nu_i} \qquad\text{(produits au numérateur, réactifs au dénominateur).} \]
Il dépend de l'état du système : il varie au cours de la transformation.
\end{definition}

\begin{rappel}[loi d'action des masses et critère d'évolution]
\begin{itemize}
  \item À l'équilibre chimique, $Q_{r,\text{éq}} = K^\circ(T)$ : la \textbf{constante thermodynamique d'équilibre} $K^\circ$ ne dépend que de la température.
  \item Si $Q_r < K^\circ$, le système évolue dans le sens direct ; si $Q_r > K^\circ$, dans le sens indirect.
\end{itemize}
En PTSI ce critère était \emph{admis}. Le chapitre~\ref{ch:evolution} le \emph{démontre} à partir du deuxième principe, et les chapitres~\ref{ch:grandeurs} et~\ref{ch:evolution} montrent comment \emph{calculer} $K^\circ$ (qui était toujours donnée en première année).
\end{rappel}

\section{Oxydoréduction}

\subsection{Couples et nombres d'oxydation}

\begin{definition}[oxydant, réducteur, demi-équation]
Un \textbf{oxydant} est une espèce capable de capter des électrons, un \textbf{réducteur} une espèce capable d'en céder. Ils forment un \textbf{couple} Ox/Red relié par la demi-équation électronique
\[ \alpha\,\mathrm{Ox} + \gamma\,e^- \;=\; \beta\,\mathrm{Red} \]
(équilibrée si besoin avec \ce{H+} et \ce{H2O}). Une \textbf{réaction d'oxydoréduction} échange $\gamma$ électrons entre l'oxydant d'un couple et le réducteur d'un autre.
\end{definition}

\begin{definition}[nombre d'oxydation]
Le \textbf{nombre d'oxydation} (n.o.) d'un élément dans un édifice est la charge fictive qu'il porterait si l'on attribuait chaque doublet de liaison à l'atome le plus électronégatif. On le note en chiffres romains. Règles pratiques :
\begin{itemize}
  \item la somme des n.o. des atomes est égale à la charge de l'édifice ;
  \item O vaut en général $-\mathrm{II}$ (sauf peroxydes : $-\mathrm{I}$) ; H vaut en général $+\mathrm{I}$ ;
  \item un corps simple (\ce{O2}, \ce{Fe}, \ce{H2}\dots) a tous ses n.o. nuls.
\end{itemize}
Une \emph{oxydation} fait \emph{croître} le n.o. ; une \emph{réduction} le fait \emph{décroître}.
\end{definition}

\begin{exemple}[calculs de n.o.]
\begin{itemize}
  \item \ce{ClO4^-} : $x + 4\times(-2) = -1$ donc $x = +\mathrm{VII}$ pour le chlore (on retrouve le résultat en attribuant tous les doublets Cl--O à l'oxygène dans la formule de Lewis).
  \item \ce{H2O2} : $2\times(+1) + 2x = 0$ donc O est au n.o. $-\mathrm{I}$ (pont peroxyde \ce{O-O}).
  \item \ce{O2^-} (ion superoxyde) : $2x = -1$ donc $x = -\tfrac12$. Un n.o. non entier est une \emph{moyenne} : en réalité, les deux atomes ne sont pas dans le même environnement électronique.
\end{itemize}
\end{exemple}

\subsection{Potentiel d'électrode et formule de Nernst}

\begin{definition}[électrode, potentiel d'électrode]
Une \textbf{électrode} (ou demi-pile) est l'ensemble formé d'un conducteur électronique (métal, graphite\dots) au contact d'un électrolyte contenant les deux membres d'un couple. Son \textbf{potentiel} $E$ est la tension entre ce conducteur et une électrode de référence : l'\textbf{électrode standard à hydrogène} (ESH), dont le potentiel est nul par convention à toute température.
\end{definition}

\begin{resultat}{Formule de Nernst}
Pour le couple $\alpha\,\mathrm{Ox} + \gamma e^- + \dots = \beta\,\mathrm{Red} + \dots$ à l'équilibre électrochimique :
\[ E = E^\circ(\mathrm{Ox/Red}) + \frac{RT}{\gamma F}\ln\frac{a_{\mathrm{Ox}}^{\,\alpha}\,(\dots)}{a_{\mathrm{Red}}^{\,\beta}\,(\dots)} \;\underset{\SI{298}{K}}{\approx}\; E^\circ + \frac{0{,}06}{\gamma}\log\frac{a_{\mathrm{Ox}}^{\,\alpha}\,(\dots)}{a_{\mathrm{Red}}^{\,\beta}\,(\dots)} \]
où $E^\circ$ est le \textbf{potentiel standard} du couple (il ne dépend que de $T$) et où les « $(\dots)$ » désignent les activités des autres espèces de la demi-équation (\ce{H+}, \ce{Cl-}\dots), chacune à la puissance de son coefficient.
\end{resultat}

\begin{exemple}[quelques électrodes classiques]
\begin{itemize}
  \item \textbf{Métal dans une solution de ses ions} : \ce{Cu^{2+} + 2e^- = Cu(s)}, donc $E = 0{,}34 + 0{,}03\log\dfrac{[\ce{Cu^{2+}}]}{c^\circ}$.
  \item \textbf{Électrode à gaz} : \ce{2H+ + 2e^- = H2(g)} sur platine, donc $E = 0 + 0{,}03\log\dfrac{([\ce{H+}]/c^\circ)^2}{P_{\ce{H2}}/P^\circ}$. Dans les conditions standard ($\mathrm{pH} = 0$, $P_{\ce{H2}} = \SI{1}{bar}$), $E = \SI{0}{V}$ : c'est l'ESH.
  \item \textbf{Électrode de deuxième espèce} (métal, sel peu soluble de ce métal, anion commun) : l'électrode au calomel \ce{Hg2Cl2(s)/Hg(l)} dans KCl, de potentiel $E = 0{,}27 - 0{,}06\log\dfrac{[\ce{Cl-}]}{c^\circ}$. Dans KCl saturé, c'est l'électrode au calomel saturée (ECS), $E_{\text{ECS}} \approx \SI{0,24}{V}$. Le calcul du « $0{,}27$ » est fait au chapitre~\ref{ch:redox}.
\end{itemize}
\end{exemple}

\subsection{Prévision des réactions : règle du gamma}

\begin{figure}[H]
\centering
\begin{tikzpicture}[>=Stealth]
  \draw[->,line width=0.9pt] (0,-0.3) -- (0,4.5) node[above] {$E^\circ$};
  \draw (-0.12,3.2) -- (0.12,3.2); \draw (-0.12,0.8) -- (0.12,0.8);
  \node[left] at (-0.25,3.2) {$\mathrm{Ox}_1$}; \node[right] at (0.25,3.2) {$\mathrm{Red}_1$};
  \node[left] at (-0.25,0.8) {$\mathrm{Ox}_2$}; \node[right] at (0.25,0.8) {$\mathrm{Red}_2$};
  \node[left,font=\small,color=black!60] at (-1.4,3.2) {$E_1^\circ$};
  \node[left,font=\small,color=black!60] at (-1.4,0.8) {$E_2^\circ$};
  \draw[c3,line width=1.2pt,->] (-0.75,3.45) to[out=30,in=150] node[above,font=\scriptsize,fill=white,inner sep=1pt] {réduction} (0.85,3.45);
  \draw[c4,line width=1.2pt,->] (0.85,0.55) to[out=-150,in=-30] node[below,font=\scriptsize,fill=white,inner sep=1pt] {oxydation} (-0.75,0.55);
  \draw[navy,line width=1pt,dashed] (-0.7,2.95) -- (0.75,1.05);
  \node[c3,font=\small,align=left,anchor=west] at (2.0,2.0) {réaction spontanée (sens direct, $K^\circ>1$) :\\[0.2em]$\mathrm{Ox}_1 + \mathrm{Red}_2 \longrightarrow \mathrm{Red}_1 + \mathrm{Ox}_2$};
\end{tikzpicture}
\caption{Règle du gamma : l'oxydant le plus fort (en haut à gauche) réagit sur le réducteur le plus fort (en bas à droite).}
\label{fig:gamma}
\end{figure}

\begin{resultat}{Constante d'équilibre d'une réaction d'oxydoréduction (PTSI)}
Pour $a_1\mathrm{Ox}_1 + b_2\mathrm{Red}_2 = b_1\mathrm{Red}_1 + a_2\mathrm{Ox}_2$ échangeant $\gamma$ électrons :
\[ \ln K^\circ = \frac{\gamma F}{RT}\left(E_1^\circ - E_2^\circ\right) \qquad\Longrightarrow\qquad \boxed{\log K^\circ \underset{\SI{298}{K}}{=} \frac{\gamma\,(E_1^\circ - E_2^\circ)}{0{,}06}} \]
\end{resultat}

\begin{demo}
À l'équilibre, la solution n'a qu'un seul potentiel : $E_1 = E_2$. En écrivant Nernst pour les deux couples (avec le même nombre $\gamma$ d'électrons, après multiplication des demi-équations) :
\[ E_1^\circ + \frac{RT}{\gamma F}\ln\frac{a_{\mathrm{Ox}_1}^{a_1}}{a_{\mathrm{Red}_1}^{b_1}} = E_2^\circ + \frac{RT}{\gamma F}\ln\frac{a_{\mathrm{Ox}_2}^{a_2}}{a_{\mathrm{Red}_2}^{b_2}} \;\Longrightarrow\; E_1^\circ - E_2^\circ = \frac{RT}{\gamma F}\ln\underbrace{\frac{a_{\mathrm{Red}_1}^{b_1}\,a_{\mathrm{Ox}_2}^{a_2}}{a_{\mathrm{Ox}_1}^{a_1}\,a_{\mathrm{Red}_2}^{b_2}}}_{Q_{r,\text{éq}} = K^\circ} \]
\end{demo}

Le chapitre~\ref{ch:redox} retrouve ce résultat par une voie purement thermodynamique, en reliant $\Delta_rG$ aux potentiels.

% ================================================================== Chapitre 2
\chapter{Enthalpie libre et potentiel chimique}
\label{ch:potentiel}

En thermodynamique physique, les transformations étudiées ont lieu le plus souvent à température et pression fixées (réacteur ouvert à l'air, bain thermostaté). La fonction d'état adaptée à ces conditions est l'\emph{enthalpie libre} $G$. Pour un système dont la composition change, on lui associe une grandeur par constituant : le \emph{potentiel chimique}.

\section{La fonction enthalpie libre}

\begin{definition}[enthalpie libre]
L'\textbf{enthalpie libre} (ou énergie de Gibbs) d'un système est la fonction d'état
\[ G = H - TS = U + PV - TS \]
où $U$ est l'énergie interne, $H = U + PV$ l'enthalpie et $S$ l'entropie. C'est une grandeur \emph{extensive}, en joules.
\end{definition}

\begin{resultat}{Identité thermodynamique -- système fermé de composition fixée}
\[ \dd G = V\,\dd P - S\,\dd T \]
Les variables naturelles de $G$ sont $T$ et $P$ : à $T$ et $P$ fixées, $G$ d'un système de composition fixée ne varie pas.
\end{resultat}

\begin{demo}
On différentie $G = U + PV - TS$ : $\dd G = \dd U + P\,\dd V + V\,\dd P - T\,\dd S - S\,\dd T$. Pour une transformation réversible élémentaire, $\dd U = \delta W + \delta Q = -P\,\dd V + T\,\dd S$, d'où
\[ \dd G = (-P\,\dd V + T\,\dd S) + P\,\dd V + V\,\dd P - T\,\dd S - S\,\dd T = V\,\dd P - S\,\dd T. \]
Comme $G$ est une fonction d'état, cette relation entre états voisins reste vraie pour toute transformation.
\end{demo}

\begin{pieges}
\begin{itemize}
  \item $G = U \mathbf{+} PV - TS$ : le signe devant $PV$ est \emph{plus} (car $H = U + PV$). Une erreur de signe ici rend fausse toute la suite du calcul de $\dd G$.
  \item $\dd G = V\,\dd P - S\,\dd T$ suppose le système \emph{fermé} et \emph{de composition fixée}. Dès qu'une réaction a lieu, il faut ajouter les termes $\mu_i\,\dd n_i$ (section suivante).
\end{itemize}
\end{pieges}

\section{Système de composition variable : potentiel chimique}

Dès qu'une transformation chimique a lieu, les quantités de matière $n_i$ varient : $G$ devient une fonction $G(T, P, n_1, n_2, \dots)$.

\begin{definition}[potentiel chimique]
Le \textbf{potentiel chimique} du constituant $\mathrm{A}_i$ dans un mélange est
\[ \mu_i = \left(\frac{\partial G}{\partial n_i}\right)_{T,\,P,\,n_{j\neq i}} \qquad\text{en }\mathrm{J\cdot mol^{-1}}. \]
C'est une grandeur \emph{intensive} : la variation d'enthalpie libre du système quand on y ajoute une mole de $\mathrm{A}_i$ en gardant $T$, $P$ et les autres quantités constantes (à la limite d'un ajout infinitésimal).
\end{definition}

\begin{resultat}{Différentielle et expression de $G$ pour un mélange}
\[ \dd G = V\,\dd P - S\,\dd T + \sum_i \mu_i\,\dd n_i \qquad\qquad G = \sum_i n_i\,\mu_i \]
\end{resultat}

La seconde relation (\emph{identité d'Euler}) est admise : elle traduit l'extensivité de $G$. On peut s'en convaincre en « remplissant » le système à $T$, $P$ et composition fixées (donc à $\mu_i$ constants) : en ajoutant les constituants dans les bonnes proportions, chaque $\dd n_i$ apporte $\mu_i\,\dd n_i$ et l'on obtient $G = \sum_i \mu_i n_i$.

\begin{exemple}[corps pur seul]
Pour un corps pur, $G = n\,\mu$, donc $\mu = G/n = G_m$ : le potentiel chimique d'un corps pur est son \textbf{enthalpie libre molaire}. Pour une mole, l'identité thermodynamique devient $\dd\mu = V_m\,\dd P - S_m\,\dd T$ avec $G_m = H_m - TS_m$.
\end{exemple}

\section{Expression du potentiel chimique}

\begin{resultat}{Potentiel chimique en fonction de l'activité (admis)}
\[ \boxed{\mu_i(T, P, \text{composition}) = \mu_i^\circ(T) + RT\ln a_i} \]
$\mu_i^\circ(T)$ est le \textbf{potentiel chimique standard} de $\mathrm{A}_i$ : sa valeur quand $a_i = 1$, c'est-à-dire quand $\mathrm{A}_i$ est dans son état standard. Il ne dépend que de $T$.
\end{resultat}

Avec le tableau des activités du chapitre précédent :
\[ \mu_i^{\text{gaz}} = \mu_i^\circ + RT\ln\frac{P_i}{P^\circ}, \qquad \mu_i^{\text{soluté}} = \mu_i^\circ + RT\ln\frac{c_i}{c^\circ}, \qquad \mu^{\text{solide ou liquide pur}} = \mu^\circ. \]

\begin{horsprog}[origine du terme $RT\ln a_i$]
Pour un gaz parfait pur, $\dd\mu = V_m\,\dd P = \frac{RT}{P}\,\dd P$ à $T$ fixée ; en intégrant de $P^\circ$ à $P$, on obtient $\mu = \mu^\circ + RT\ln(P/P^\circ)$. Les autres expressions s'obtiennent par analogie (modèle de la solution idéale). Seul le résultat est au programme.
\end{horsprog}

\section{États standard}

\begin{definition}[état standard d'un constituant à la température $T$]
L'\textbf{état standard} d'un constituant est un état de référence, choisi par convention à la \emph{même température $T$} que le système et sous la \textbf{pression standard} $P^\circ = \SI{1}{bar}$ :
\begin{itemize}
  \item \textbf{gaz} : le constituant pur, se comportant comme un gaz parfait, sous $P^\circ$ ;
  \item \textbf{solide ou liquide} : le constituant pur, dans le même état physique, sous $P^\circ$ ;
  \item \textbf{soluté} : le soluté à la concentration $c^\circ = \SI{1}{mol.L^{-1}}$, se comportant comme en solution infiniment diluée ;
  \item \textbf{solvant} : le solvant pur, sous $P^\circ$.
\end{itemize}
\end{definition}

\begin{pieges}
\begin{itemize}
  \item Il n'y a \emph{pas} de température standard : il existe un état standard \emph{à chaque température}. Les tables sont souvent données à $\SI{298}{K}$ par commodité, mais « standard » ne veut pas dire « $\SI{25}{\celsius}$ ».
  \item L'état standard peut être \emph{fictif} : un gaz réel pur n'est jamais exactement parfait ; un soluté à $\SI{1}{mol.L^{-1}}$ n'est jamais infiniment dilué. C'est une référence de calcul, pas un état réel.
  \item Pour \ce{H+} en solution, l'état standard correspond à $[\ce{H+}] = c^\circ$, donc à $\mathrm{pH} = 0$ : c'est pourquoi les potentiels standard des couples faisant intervenir \ce{H+} sont ceux mesurés à $\mathrm{pH} = 0$.
\end{itemize}
\end{pieges}

\begin{definition}[état standard de référence d'un élément]
L'\textbf{état standard de référence} (ESR) d'un élément chimique à la température $T$ est l'état standard du \emph{corps simple} correspondant, dans l'état physique \emph{le plus stable} à cette température (sous $P^\circ$).

\noindent\emph{Exceptions conventionnelles} : pour H, N, O, F et Cl, l'ESR est le gaz diatomique (\ce{H2}, \ce{N2}, \ce{O2}, \ce{F2}, \ce{Cl2}) \emph{à toute température} ; pour le carbone, c'est le graphite.
\end{definition}

\begin{exemple}[ESR à différentes températures]
\begin{itemize}
  \item Oxygène : \ce{O2(g)} et non \ce{O3(g)} (l'ozone est un corps simple moins stable).
  \item Brome : \ce{Br2(l)} à $\SI{298}{K}$, mais \ce{Br2(g)} au-dessus de sa température d'ébullition $\SI{332}{K}$.
  \item Fer : \ce{Fe(s)} jusqu'à $\SI{1811}{K}$, puis \ce{Fe(l)}.
  \item Carbone : graphite et non diamant.
\end{itemize}
\end{exemple}

\begin{exemple}[graphite ou diamant ?]
À $\SI{298}{K}$, $\mu^\circ(\text{C}_{\text{diamant}}) - \mu^\circ(\text{C}_{\text{graphite}}) = +\SI{2,9}{kJ.mol^{-1}} > 0$. Les deux étant des solides purs ($a = 1$), $\mu = \mu^\circ$ : le diamant a un potentiel chimique plus élevé. La transformation diamant $\to$ graphite est donc thermodynamiquement favorisée (on le démontre au chapitre~\ref{ch:evolution}), mais elle est infiniment lente à température ambiante : le diamant est \emph{métastable}. La thermodynamique dit ce qui \emph{peut} se produire, la cinétique dit \emph{à quelle vitesse}.
\end{exemple}

% ================================================================== Chapitre 3
\chapter{Grandeurs de réaction et grandeurs standard}
\label{ch:grandeurs}

On considère un système siège d'\emph{une seule} réaction $0 = \sum_i\nu_i\mathrm{A}_i$, d'avancement $\xi$. Toute grandeur extensive $X$ du système ($G$, $H$, $S$, $V$\dots) dépend alors de $T$, $P$ et $\xi$. La question centrale de ce chapitre : \emph{comment $X$ varie-t-elle quand la réaction avance ?}

\section{Grandeurs de réaction}

\begin{definition}[grandeur molaire partielle]
La \textbf{grandeur molaire partielle} de $\mathrm{A}_i$ associée à $X$ est $X_{m,i} = \left(\dfrac{\partial X}{\partial n_i}\right)_{T,P,n_{j\neq i}}$. Pour $X = G$, c'est le potentiel chimique : $G_{m,i} = \mu_i$.
\end{definition}

\begin{definition}[grandeur de réaction]
La \textbf{grandeur de réaction} associée à $X$ est la dérivée de $X$ par rapport à l'avancement, à $T$ et $P$ fixées :
\[ \boxed{\Delta_r X = \left(\frac{\partial X}{\partial \xi}\right)_{T,P} = \sum_i \nu_i\,X_{m,i}} \qquad\text{en particulier}\qquad \boxed{\Delta_r G = \sum_i \nu_i\,\mu_i} \]
Elle s'exprime en $\mathrm{J\cdot mol^{-1}}$ (ou $\mathrm{J\cdot K^{-1}\cdot mol^{-1}}$ pour $\Delta_rS$) : « par mole d'avancement ». Elle dépend de $T$, de $P$ \emph{et de la composition} du système.
\end{definition}

\begin{demo}
À $T$ et $P$ fixées, $\dd X = \sum_i X_{m,i}\,\dd n_i = \sum_i X_{m,i}\,\nu_i\,\dd\xi$, d'où $\left(\partial X/\partial\xi\right)_{T,P} = \sum_i\nu_iX_{m,i}$.
\end{demo}

Le symbole $\Delta_r$ n'est \emph{pas} une différence entre deux états : c'est un opérateur de dérivation (« opérateur de Lewis »). Une grandeur de réaction est une \emph{pente} : $\Delta_rG$ est la pente de la courbe $G(\xi)$ (chapitre~\ref{ch:evolution}).

\begin{pieges}
\begin{itemize}
  \item Une grandeur de réaction est attachée à une \emph{équation écrite} : si l'on multiplie tous les coefficients par 2, $\Delta_rX$ est multipliée par 2. Toujours préciser l'équation utilisée.
  \item Ne pas confondre $\Delta_rH$ (en $\mathrm{J\cdot mol^{-1}}$, intensive) et $\Delta H$ (en J, variation entre deux états). Le lien est $\Delta H = \int \Delta_rH\,\dd\xi$.
\end{itemize}
\end{pieges}

\section{Grandeurs standard de réaction}

\begin{definition}[grandeur standard de réaction]
La \textbf{grandeur standard de réaction} est la grandeur de réaction calculée en supposant \emph{chaque constituant dans son état standard} à la température $T$ :
\[ \boxed{\Delta_r X^\circ(T) = \sum_i \nu_i\,X^\circ_{m,i}(T)} \]
en particulier
\[ \Delta_rG^\circ = \sum_i\nu_i\,\mu_i^\circ,\qquad \Delta_rH^\circ = \sum_i\nu_iH^\circ_{m,i},\qquad \Delta_rS^\circ = \sum_i\nu_iS^\circ_{m,i}. \]
Les $X^\circ_{m,i}$ sont les grandeurs molaires des constituants purs dans leur état standard : $\Delta_rX^\circ$ \emph{ne dépend que de $T$}.
\end{definition}

On peut se représenter $\Delta_rX^\circ$ comme la variation de $X$ pour une transformation \emph{fictive} : une mole d'avancement, à $T$ et $P^\circ$ fixées, où réactifs et produits resteraient chacun dans leur état standard (donc purs et séparés). Une telle réaction « totale fictive » n'est pas réalisable, mais elle définit une référence tabulable.

\begin{figure}[H]
\centering
\begin{tikzpicture}[>=Stealth,box/.style={draw=navy,line width=0.8pt,rounded corners=2pt,fill=cdefbg,minimum width=3.3cm,minimum height=1cm,align=center,font=\small}]
  \node[box] (R) at (0,0) {Réactifs\\dans l'état standard};
  \node[box] (P) at (7.4,0) {Produits\\dans l'état standard};
  \draw[->,line width=1.1pt,c3] (R) -- node[above,font=\small,black] {réaction fictive, totale} node[below,font=\small,black] {$T$ fixée, $P = P^\circ$} (P);
  \node[font=\small,c3] at (3.7,-1.05) {$\Delta X = \Delta_rX^\circ$ pour $\xi = \SI{1}{mol}$};
\end{tikzpicture}
\caption{Interprétation d'une grandeur standard de réaction.}
\end{figure}

\begin{resultat}{Relation fondamentale}
\[ \boxed{\Delta_rG^\circ(T) = \Delta_rH^\circ(T) - T\,\Delta_rS^\circ(T)} \]
\end{resultat}

\begin{demo}
Pour chaque constituant dans son état standard, $\mu_i^\circ = G^\circ_{m,i} = H^\circ_{m,i} - TS^\circ_{m,i}$. On multiplie par $\nu_i$ et on somme.
\end{demo}

\section{Signification physique et signe}

\begin{definition}[réaction endothermique, exothermique, athermique]
Une réaction est dite \textbf{exothermique} si $\Delta_rH^\circ < 0$, \textbf{endothermique} si $\Delta_rH^\circ > 0$, \textbf{athermique} si $\Delta_rH^\circ \approx 0$. Le sens physique (le système cède ou reçoit de la chaleur) est établi au chapitre~\ref{ch:premier}.
\end{definition}

\begin{resultat}{Signe de l'entropie standard de réaction}
L'entropie molaire d'un gaz (de l'ordre de $\SI{200}{J.K^{-1}.mol^{-1}}$) est très supérieure à celle d'un solide ou d'un liquide (quelques dizaines). En notant $\Delta\nu_{\text{gaz}} = \sum_{i\,\text{gaz}}\nu_i$ :
\[ \Delta\nu_{\text{gaz}} > 0 \;\Rightarrow\; \Delta_rS^\circ > 0, \qquad \Delta\nu_{\text{gaz}} < 0 \;\Rightarrow\; \Delta_rS^\circ < 0, \qquad \Delta\nu_{\text{gaz}} = 0 \;\Rightarrow\; |\Delta_rS^\circ| \text{ faible}. \]
Interprétation : l'entropie mesure le désordre ; créer du gaz augmente fortement le nombre de configurations accessibles.
\end{resultat}

\begin{exemple}[synthèse de l'ammoniac]
\ce{N2(g) + 3 H2(g) = 2 NH3(g)} : $\Delta\nu_{\text{gaz}} = 2 - 4 = -2 < 0$, donc on prévoit $\Delta_rS^\circ < 0$. Avec les entropies molaires standard (tableau~\ref{tab:donnees}) :
\begin{align*} \Delta_rS^\circ &= 2S^\circ_m(\ce{NH3}) - S^\circ_m(\ce{N2}) - 3S^\circ_m(\ce{H2})\\ &= 2(192{,}8) - 191{,}6 - 3(130{,}7) = \SI{-198,1}{J.K^{-1}.mol^{-1}}. \end{align*}
\end{exemple}

\section{Grandeurs standard de formation et loi de Hess}

\begin{definition}[réaction de formation, grandeur standard de formation]
La \textbf{réaction standard de formation} d'un constituant à la température $T$ est la réaction qui forme \emph{une mole} de ce constituant, dans son état standard, à partir des éléments qui le composent pris dans leur \emph{état standard de référence}.

\noindent La \textbf{grandeur standard de formation} $\Delta_fX^\circ$ est la grandeur standard de cette réaction.
\end{definition}

\begin{exemple}
Formation de l'ammoniac : $\tfrac12\,\ce{N2(g)} + \tfrac32\,\ce{H2(g)} \longrightarrow \ce{NH3(g)}$, avec $\Delta_fH^\circ(\ce{NH3(g)}) = \SI{-45,9}{kJ.mol^{-1}}$ (donnée de table).

\noindent Formation de l'ozone : $\tfrac32\,\ce{O2(g)} \longrightarrow \ce{O3(g)}$, donc $\Delta_fX^\circ(\ce{O3}) \neq 0$ alors que c'est un corps simple.
\end{exemple}

\begin{resultat}{Conséquence immédiate}
\[ \boxed{\Delta_fX^\circ(\text{corps simple dans son ESR}) = 0} \]
car sa « réaction de formation » ne transforme rien (\ce{O2(g) -> O2(g)}). En revanche $\Delta_fX^\circ \neq 0$ pour un corps composé \emph{ou} pour un corps simple hors de son ESR (\ce{O3}, diamant, \ce{Br2(g)} à $\SI{298}{K}$\dots).
\end{resultat}

\begin{resultat}{Loi de Hess}
\[ \boxed{\Delta_rH^\circ = \sum_i \nu_i\,\Delta_fH^\circ_i} \qquad\text{et de même}\qquad \Delta_rG^\circ = \sum_i\nu_i\,\Delta_fG^\circ_i \]
\end{resultat}

\begin{demo}
$H$ est une fonction d'état : sa variation ne dépend pas du chemin suivi. On passe des réactifs aux produits en deux étapes fictives : (1) décomposer les réactifs en éléments dans leur ESR, ce qui contribue $-\sum_{\text{réactifs}}|\nu_i|\,\Delta_fH^\circ_i$ ; (2) reformer les produits à partir de ces éléments, ce qui contribue $+\sum_{\text{produits}}\nu_i\,\Delta_fH^\circ_i$. Avec les coefficients algébriques, les deux sommes n'en font qu'une.
\end{demo}

\begin{figure}[H]
\centering
\begin{tikzpicture}[>=Stealth,box/.style={draw=navy,line width=0.8pt,rounded corners=2pt,fill=cdefbg,inner sep=5pt,align=center,font=\small}]
  \node[box] (R) at (0,2.6) {\ce{CH4(g) + Cl2(g)}};
  \node[box] (P) at (8.6,2.6) {\ce{CH3Cl(g) + HCl(g)}};
  \node[box,fill=cmetbg,draw=c4] (E) at (4.3,0) {\ce{C(graphite) + 2 H2(g) + Cl2(g)}\\{\scriptsize éléments dans leur ESR}};
  \draw[->,line width=1.2pt,c3] (R) -- node[above,font=\small] {$\Delta_rH^\circ = \;?$} (P);
  \draw[->,line width=1pt,navy] (R.south) -- node[below left,font=\small,align=right,pos=0.5] {$-\Delta_fH^\circ(\ce{CH4})$\\$-\,\Delta_fH^\circ(\ce{Cl2}) = 0$} (E.north west);
  \draw[->,line width=1pt,navy] (E.north east) -- node[below right,font=\small,align=left,pos=0.5] {$\Delta_fH^\circ(\ce{CH3Cl})$\\$+\,\Delta_fH^\circ(\ce{HCl})$} (P.south);
\end{tikzpicture}
\caption{Cycle de Hess pour la chloration du méthane : on passe par les éléments dans leur état standard de référence.}
\label{fig:hess}
\end{figure}

\begin{exemple}[chloration du méthane]
\ce{CH4(g) + Cl2(g) -> CH3Cl(g) + HCl(g)}. D'après la loi de Hess et les données à $\SI{298}{K}$ :
\begin{align*} \Delta_rH^\circ &= \Delta_fH^\circ(\ce{CH3Cl}) + \Delta_fH^\circ(\ce{HCl}) - \Delta_fH^\circ(\ce{CH4}) - \underbrace{\Delta_fH^\circ(\ce{Cl2})}_{0}\\ &= -81{,}9 - 92{,}3 + 74{,}6 = \SI{-99,6}{kJ.mol^{-1}}. \end{align*}
La réaction est exothermique.
\end{exemple}

\begin{resultat}{Et pour l'entropie ?}
Les tables ne donnent pas $\Delta_fS^\circ$ mais les \textbf{entropies molaires standard absolues} $S^\circ_m$ (le troisième principe fixe $S_m = 0$ pour un cristal parfait à $\SI{0}{K}$). On calcule donc directement
\[ \Delta_rS^\circ = \sum_i\nu_i\,S^\circ_{m,i}. \]
\end{resultat}

\begin{pieges}
\begin{itemize}
  \item $\Delta_fH^\circ(\ce{O2}) = 0$ \emph{mais} $S^\circ_m(\ce{O2}) = \SI{205,2}{J.K^{-1}.mol^{-1}} \neq 0$ : les entropies des corps simples ne sont \emph{jamais} nulles. Les oublier dans $\Delta_rS^\circ$ est l'erreur la plus fréquente.
  \item Unités : $\Delta_rH^\circ$ en $\mathrm{kJ\cdot mol^{-1}}$, $\Delta_rS^\circ$ en $\mathrm{J\cdot K^{-1}\cdot mol^{-1}}$. Dans $\Delta_rG^\circ = \Delta_rH^\circ - T\Delta_rS^\circ$, convertir avant de soustraire.
  \item Bien vérifier l'état physique : $\Delta_fH^\circ(\ce{H2O(l)}) = -285{,}8$ et $\Delta_fH^\circ(\ce{H2O(g)}) = \SI{-241,8}{kJ.mol^{-1}}$ ; la différence est l'enthalpie de vaporisation.
\end{itemize}
\end{pieges}

\section{Combinaisons linéaires de réactions}

\begin{resultat}{Combinaison linéaire}
Si l'équation (3) s'obtient comme $(3) = a\,(1) + b\,(2)$, alors
\[ \Delta_rX_3^\circ = a\,\Delta_rX_1^\circ + b\,\Delta_rX_2^\circ \quad (X = H, S, G) \qquad\text{et}\qquad K_3^\circ = \left(K_1^\circ\right)^a\left(K_2^\circ\right)^b. \]
\end{resultat}

\begin{demo}
$\Delta_rX^\circ = \sum_i\nu_iX^\circ_{m,i}$ est linéaire en les coefficients $\nu_i$, d'où la première relation. Pour $K^\circ$, on utilise $\Delta_rG^\circ = -RT\ln K^\circ$ (établi au chapitre~\ref{ch:evolution}) : $\ln K_3^\circ = a\ln K_1^\circ + b\ln K_2^\circ$.
\end{demo}

\begin{exemple}[enthalpie de formation de CO]
On ne peut pas mesurer directement \ce{C + 1/2 O2 -> CO} (il se forme toujours un peu de \ce{CO2}). On combine :
\begin{align*} (1)&\quad \ce{C(gr) + O2 -> CO2} & \Delta_rH_1^\circ &= \SI{-393,5}{kJ.mol^{-1}}\\ (2)&\quad \ce{CO + 1/2 O2 -> CO2} & \Delta_rH_2^\circ &= \SI{-283,0}{kJ.mol^{-1}} \end{align*}
$(3) = (1) - (2)$ donne \ce{C(gr) + 1/2 O2 -> CO}, d'où $\Delta_fH^\circ(\ce{CO}) = -393{,}5 + 283{,}0 = \SI{-110,5}{kJ.mol^{-1}}$.
\end{exemple}

\section{Influence de la température : approximation d'Ellingham}

\begin{definition}[approximation d'Ellingham]
On suppose que $\Delta_rH^\circ$ et $\Delta_rS^\circ$ sont \textbf{indépendantes de la température} sur tout intervalle où aucun constituant ne change d'état. Alors
\[ \boxed{\Delta_rG^\circ(T) = \Delta_rH^\circ - T\,\Delta_rS^\circ} \]
est une \textbf{fonction affine} de $T$, de pente $-\Delta_rS^\circ$ et d'ordonnée à l'origine $\Delta_rH^\circ$.
\end{definition}

Cette approximation est imposée par le programme. Elle est justifiée car $\Delta_rH^\circ$ et $\Delta_rS^\circ$ varient peu avec $T$ (leurs dérivées font intervenir $\Delta_rC_p^\circ$, petite devant $\Delta_rH^\circ/T$), alors que le terme $T\Delta_rS^\circ$ change beaucoup.

\begin{figure}[H]
\centering
\begin{tikzpicture}
\begin{axis}[cours,xlabel={$T$ (K)},ylabel={$\Delta_rG^\circ$ ($\mathrm{kJ\cdot mol^{-1}}$)},
  xmin=300,xmax=900,ymin=-40,ymax=100,
  xtick={300,400,500,600,700,800,900}]
  \fill[c4!12] (axis cs:300,0) -- (axis cs:300,-32.37) -- (axis cs:463.4,0) -- cycle;
  \fill[c3!10] (axis cs:463.4,0) -- (axis cs:900,86.49) -- (axis cs:900,0) -- cycle;
  \draw[black,line width=0.6pt] (axis cs:300,0) -- (axis cs:900,0);
  \addplot[c1] table[x=T,y=DrG_kJ]{data/nh3_drg.dat};
  \addplot[c3,only marks,mark=*,mark size=2.2pt] coordinates {(463,0)};
  \node[c3,font=\small,anchor=south east] at (axis cs:458,3) {$T_i = \SI{463}{K}$};
  \node[c4,font=\scriptsize] at (axis cs:340,-8) {$K^\circ>1$};
  \node[c3,font=\small,align=center] at (axis cs:790,28) {$\Delta_rG^\circ>0$\\$K^\circ<1$};
  \node[font=\small,anchor=west] at (axis cs:310,80) {pente $= -\Delta_rS^\circ = +\SI{0,198}{kJ.K^{-1}.mol^{-1}}$};
\end{axis}
\end{tikzpicture}
\caption{$\Delta_rG^\circ(T)$ pour \ce{N2 + 3H2 = 2NH3} dans l'approximation d'Ellingham ($\Delta_rH^\circ = \SI{-91,8}{kJ.mol^{-1}}$, $\Delta_rS^\circ = \SI{-198,1}{J.K^{-1}.mol^{-1}}$).}
\label{fig:nh3drg}
\end{figure}

\begin{definition}[température d'inversion]
La \textbf{température d'inversion} $T_i$ est la température où $\Delta_rG^\circ(T_i) = 0$, soit $K^\circ(T_i) = 1$. Dans l'approximation d'Ellingham, $T_i = \Delta_rH^\circ/\Delta_rS^\circ$.
\end{definition}

\begin{exemple}[ammoniac, suite]
$\Delta_rG^\circ(\SI{298}{K}) = -91{,}8 - 298\times(-0{,}1981) = \SI{-32,7}{kJ.mol^{-1}}$ et $T_i = \dfrac{-91{,}8\times10^3}{-198{,}1} = \SI{463}{K}$ ($\SI{190}{\celsius}$). La synthèse est favorisée ($K^\circ > 1$) en dessous de $T_i$ : le terme enthalpique favorable domine à basse température, le terme entropique défavorable domine à haute température.
\end{exemple}

\begin{ordres}
\begin{itemize}
  \item $|\Delta_rH^\circ|$ : quelques dizaines à quelques centaines de $\mathrm{kJ\cdot mol^{-1}}$ (combustions : $\sim\SI{-800}{kJ.mol^{-1}}$ pour \ce{CH4}).
  \item $|\Delta_rS^\circ|$ : de l'ordre de $\SI{100}{}$ à $\SI{200}{J.K^{-1}.mol^{-1}}$ par mole de gaz créée ou consommée.
  \item À $\SI{298}{K}$, $T|\Delta_rS^\circ| \sim 30$ à $\SI{60}{kJ.mol^{-1}}$ : le terme entropique ne l'emporte sur le terme enthalpique qu'à haute température ou pour des réactions peu énergétiques.
\end{itemize}
\end{ordres}

% ================================================================== Chapitre 4
\chapter{Premier principe : effets thermiques des réactions}
\label{ch:premier}

Une réaction chimique libère ou absorbe de l'énergie. Selon que le réacteur échange de la chaleur avec l'extérieur ou non, cette énergie se manifeste par un \emph{transfert thermique} ou par une \emph{variation de température}. L'outil est le premier principe, écrit avec l'enthalpie.

\section{Vocabulaire des transformations}

\begin{definition}[conditions opératoires]
Pour une transformation entre un état initial (I) et un état final (F) :
\begin{itemize}
  \item \textbf{monobare} : la pression extérieure est constante et $P_I = P_F = P_{\text{ext}}$ (réacteur ouvert sur l'atmosphère) ;
  \item \textbf{isobare} : la pression du système est définie et constante à chaque instant ;
  \item \textbf{monotherme} : le système n'échange de chaleur qu'avec un thermostat de température $T_{\text{ext}}$ ;
  \item \textbf{isotherme} : $T$ est définie et constante tout au long de la transformation (en pratique : $T_I = T_F = T_{\text{ext}}$ pour un réacteur thermostaté) ;
  \item \textbf{adiabatique} : aucun transfert thermique ($Q = 0$), par exemple dans un calorimètre ou pour une réaction très rapide.
\end{itemize}
\end{definition}

\begin{rappel}[premier principe en enthalpie]
Pour un système fermé, au repos macroscopique, subissant une transformation \emph{monobare} où seules les forces de pression travaillent :
\[ \Delta H = Q \]
($Q$ : transfert thermique \emph{reçu} par le système, algébrique).
\end{rappel}

\section{Variation d'enthalpie due à la réaction}

À $T$ et $P$ fixées, $\dd H = \Delta_rH\,\dd\xi$. On utilise ensuite l'approximation suivante, toujours faite au programme :
\[ \Delta_rH \simeq \Delta_rH^\circ. \]
Elle est justifiée car l'enthalpie d'un gaz parfait ne dépend que de $T$ (pas de sa pression partielle), celle d'une phase condensée dépend très peu de $P$, et le mélange idéal se fait sans effet thermique. Ainsi, pour une réaction menée de $\xi = 0$ à $\xi_F$ à température constante :
\[ \boxed{\Delta H = \xi_F\,\Delta_rH^\circ(T)} \]

\section{Réacteur monobare et isotherme : transfert thermique}

\begin{resultat}{Transfert thermique en réacteur monobare isotherme}
\[ \boxed{Q = \Delta H = \xi_F\,\Delta_rH^\circ} \]
\begin{itemize}
  \item réaction \textbf{exothermique} ($\Delta_rH^\circ<0$) : $Q<0$, le système \emph{cède} de la chaleur au milieu extérieur ;
  \item réaction \textbf{endothermique} ($\Delta_rH^\circ>0$) : $Q>0$, le système \emph{reçoit} de la chaleur.
\end{itemize}
\end{resultat}

\begin{exemple}[chauffage au gaz naturel]
On brûle $\SI{1}{m^3}$ de méthane (pris à $\SI{1}{bar}$ et $\SI{298}{K}$) : $n = \dfrac{PV}{RT} = \dfrac{10^5\times1}{8{,}314\times298} = \SI{40,3}{mol}$.

\noindent Réaction : \ce{CH4(g) + 2 O2(g) -> CO2(g) + 2 H2O(l)}, avec (loi de Hess)
\[ \Delta_rH^\circ = -393{,}5 + 2(-285{,}8) - (-74{,}6) = \SI{-890,5}{kJ.mol^{-1}}. \]
Si les produits sont ramenés à $\SI{298}{K}$ (eau condensée) : $Q = \xi_F\Delta_rH^\circ = 40{,}3\times(-890{,}5) = \SI{-35,9}{MJ} \approx \SI{-10}{kWh}$. C'est l'ordre de grandeur du « pouvoir calorifique » facturé par les fournisseurs de gaz ($\approx\SI{10}{kWh}$ par $\si{m^3}$).
\end{exemple}

\section{Réacteur monobare et adiabatique : température finale}

Si le réacteur est calorifugé (ou la réaction très rapide), $Q = 0$, donc $\Delta H = 0$ : l'énergie libérée par une réaction exothermique \emph{chauffe} le contenu du réacteur.

Pour calculer $T_F$, on remplace la transformation réelle (réaction et échauffement simultanés) par un \textbf{chemin fictif} en deux étapes, ce qui est légitime puisque $H$ est une fonction d'état.

\begin{figure}[H]
\centering
\begin{tikzpicture}[>=Stealth,box/.style={draw=navy,line width=0.8pt,rounded corners=2pt,fill=cdefbg,inner sep=5pt,align=center,font=\small,minimum width=3.6cm}]
  \node[box] (I) at (0,0) {État initial\\réactifs, $T_I$, $P^\circ$};
  \node[box] (A) at (8,0) {État intermédiaire\\système final, $T_I$, $P^\circ$};
  \node[box,fill=cresbg,draw=c3] (F) at (8,3.4) {État final\\système final, $T_F$, $P^\circ$};
  \draw[->,line width=1.2pt,c3,dashed] (I.north) -- node[above left,font=\small,align=center] {transformation réelle\\adiabatique : $\Delta H = 0$} (F.west);
  \draw[->,line width=1.1pt,navy] (I) -- node[above,font=\small] {(a) réaction à $T_I$} node[below,font=\small] {$\Delta H_a = \xi_F\,\Delta_rH^\circ(T_I)$} (A);
  \draw[->,line width=1.1pt,navy] (A) -- node[right,font=\small,align=left] {(b) échauffement :\\$\Delta H_b = \Big(\sum_j n_jC^\circ_{p,m,j}\Big)(T_F-T_I)$} (F);
\end{tikzpicture}
\caption{Chemin fictif pour le calcul d'une température finale en réacteur adiabatique. La somme porte sur \emph{tous} les constituants présents dans l'état final.}
\label{fig:cheminfictif}
\end{figure}

\begin{resultat}{Température finale en réacteur adiabatique monobare}
$\Delta H_a + \Delta H_b = 0$, d'où, sans changement d'état au cours de l'échauffement :
\[ \boxed{T_F = T_I - \frac{\xi_F\,\Delta_rH^\circ(T_I)}{\sum_j n_j\,C^\circ_{p,m,j}}} \]
\end{resultat}

\begin{methode}[température de fin de réaction en réacteur adiabatique]
\begin{enumerate}
  \item Écrire la réaction et dresser le tableau d'avancement ; trouver $\xi_F$ (souvent $\xi_{\max}$ pour une réaction totale).
  \item Faire l'inventaire de l'état final : produits, réactifs en excès \emph{et espèces spectatrices} (le diazote de l'air, le solvant\dots).
  \item Écrire $\Delta H = \Delta H_a + \Delta H_b = 0$ le long du chemin fictif.
  \item Si un constituant change d'état pendant l'échauffement, couper l'étape (b) à la température de changement d'état et ajouter $n\,\Delta_{\text{vap}}H^\circ$ (ou $\Delta_{\text{fus}}H^\circ$).
  \item Vérifier la cohérence : $T_F > T_I$ pour une réaction exothermique.
\end{enumerate}
\end{methode}

\begin{definition}[température de flamme adiabatique]
La \textbf{température de flamme} d'une combustion est la température finale atteinte par les gaz brûlés lorsque la combustion, totale, a lieu dans des conditions adiabatiques et monobares.
\end{definition}

\begin{exemple}[flamme méthane--air]
Mélange stœchiométrique dans l'air (environ 4 \ce{N2} pour 1 \ce{O2}) à $T_I = \SI{298}{K}$ : on part de \ce{CH4 + 2 O2 + 8 N2}, et l'état final est \ce{CO2 + 2 H2O(g) + 8 N2}. On choisit dans l'étape (a) de produire directement de l'eau \emph{gazeuse} (état standard gaz à $\SI{298}{K}$, fictif mais légitime), ce qui évite de traiter sa vaporisation :
\begin{gather*} \Delta_rH^\circ = -393{,}5 + 2(-241{,}8) + 74{,}6 = \SI{-802,5}{kJ.mol^{-1}},\\ \sum_jn_jC^\circ_{p,m,j} = 37{,}1 + 2(33{,}6) + 8(29{,}1) = \SI{337,1}{J.K^{-1}}. \end{gather*}
\[ T_F = 298 + \frac{802{,}5\times10^3}{337{,}1} \approx \SI{2680}{K}. \]
La valeur mesurée est plus proche de $\SI{2200}{K}$ : les capacités thermiques augmentent avec $T$, une partie de \ce{CO2} et de \ce{H2O} se dissocie (réactions endothermiques) et il y a des pertes. Le modèle donne néanmoins le bon ordre de grandeur. Le diazote, qui « ne fait rien », absorbe près de 70\,\% de l'énergie : c'est pourquoi une flamme dans \ce{O2} pur est beaucoup plus chaude (exercice~\ref{exo:gaz}).
\end{exemple}

\section{Capacité numérique : simulation d'un réacteur adiabatique}

Si la vitesse de la réaction est connue, on peut suivre $T(t)$. La vitesse augmente avec la température (loi d'Arrhenius $k(T) = A\,\mathrm{e}^{-E_a/RT}$) : une réaction exothermique en réacteur adiabatique \emph{s'auto-accélère}.

\begin{numerique}[modèle]
Réaction \ce{A -> B} d'ordre 1, dans un réacteur adiabatique monobare de capacité thermique totale $C$ :
\[ \frac{\dd\xi}{\dd t} = k(T)\,(n_0-\xi), \qquad C\,\dd T + \Delta_rH^\circ\,\dd\xi = 0 \;\Longrightarrow\; \frac{\dd T}{\dd t} = -\frac{\Delta_rH^\circ}{C}\,\frac{\dd\xi}{\dd t}. \]
Le bilan d'enthalpie entre $t$ et $t+\dd t$ est exactement celui du chemin fictif, appliqué à chaque pas de temps. On résout ce système par la méthode d'Euler explicite.
\end{numerique}

\begin{tcolorbox}[enhanced,breakable,colback=black!3,colframe=black!25,boxrule=0.4pt,arc=1mm,left=2mm,right=2mm]
\lstinputlisting{code/reacteur_adiabatique.py}
\end{tcolorbox}

\begin{figure}[H]
\centering
\begin{tikzpicture}
\begin{axis}[cours,width=0.78\linewidth,height=0.42\linewidth,xlabel={$t$ (min)},ylabel={$T$ ($\si{\celsius}$)},
  xmin=0,xmax=60,ymin=20,ymax=70,axis y line*=left,
  legend style={at={(0.03,0.97)},anchor=north west}]
  \addplot[c3] table[x=tmin,y=TC]{data/adiab.dat}; \addlegendentry{$T(t)$, adiabatique}
  \addplot[c3,dashed,line width=0.8pt] coordinates {(0,65) (60,65)};
  \node[c3,font=\small,anchor=south east] at (axis cs:59,65) {$T_I + \Delta T_{\text{ad}} = \SI{65}{\celsius}$};
\end{axis}
\begin{axis}[cours,width=0.78\linewidth,height=0.42\linewidth,xmin=0,xmax=60,ymin=0,ymax=1.25,
  axis x line=none,axis y line*=right,ylabel={$\xi$ (mol)},grid=none,
  legend style={at={(0.97,0.03)},anchor=south east}]
  \addplot[c1] table[x=tmin,y=xi]{data/adiab.dat}; \addlegendentry{$\xi(t)$, adiabatique}
  \addplot[c1,dashed] table[x=tmin,y=xiiso]{data/adiab.dat}; \addlegendentry{$\xi(t)$, isotherme $\SI{25}{\celsius}$}
\end{axis}
\end{tikzpicture}
\caption{Simulation ($n_0 = \SI{1}{mol}$, $\Delta_rH^\circ = \SI{-120}{kJ.mol^{-1}}$, $C = \SI{3,0}{kJ.K^{-1}}$, $E_a = \SI{80}{kJ.mol^{-1}}$, $k(\SI{298}{K}) = \SI{2,0e-4}{s^{-1}}$). Le temps de demi-réaction passe de $\SI{58}{min}$ (isotherme) à $\SI{22}{min}$ (adiabatique).}
\label{fig:adiab}
\end{figure}

\noindent\textbf{Lecture de la simulation.} La température finale vaut $T_I + \Delta T_{\text{ad}}$ avec $\Delta T_{\text{ad}} = -n_0\Delta_rH^\circ/C = \SI{40}{K}$ : on retrouve la formule du chemin fictif. Entre-temps, l'échauffement multiplie la constante de vitesse par
\[ \frac{k(T_F)}{k(T_I)} = \exp\left[\frac{E_a}{R}\left(\frac{1}{T_I}-\frac{1}{T_F}\right)\right] \approx 45, \]
d'où l'allure en S de la courbe $T(t)$, caractéristique de l'\emph{emballement thermique}. C'est le phénomène que l'on cherche à éviter en génie chimique (refroidissement, dilution, ajout progressif des réactifs).

% ================================================================== Chapitre 5
\chapter{Deuxième principe : évolution et équilibre}
\label{ch:evolution}

On établit ici ce qui était admis en première année : le critère d'évolution d'un système chimique et la loi d'action des masses. Tout découle d'une seule idée : à $T$ et $P$ fixées, l'enthalpie libre d'un système ne peut que diminuer.

\noindent\emph{Cadre de tout le chapitre} : système fermé, siège d'une réaction unique, maintenu à température et pression fixées ($T = T_{\text{ext}}$, $P = P_{\text{ext}}$), où seules les forces de pression travaillent.

\section{Création d'entropie et enthalpie libre de réaction}

\begin{resultat}{Création d'entropie due à la réaction}
\[ \boxed{\delta S_c = -\frac{\Delta_rG\,\dd\xi}{T} \geq 0} \qquad\Longleftrightarrow\qquad \dd G = -T\,\delta S_c \leq 0 \]
\end{resultat}

\begin{demo}
Sur une évolution élémentaire :
\begin{itemize}
  \item premier principe (monobare, forces de pression seules) : $\dd H = \delta Q$ ;
  \item deuxième principe (monotherme, à la température $T$ du thermostat) : $\dd S = \dfrac{\delta Q}{T} + \delta S_c$ avec $\delta S_c \geq 0$ ;
  \item $T$ étant fixée : $\dd G = \dd H - T\,\dd S = \delta Q - \delta Q - T\,\delta S_c = -T\,\delta S_c$.
\end{itemize}
Par ailleurs, à $T$ et $P$ fixées, $\dd G = \sum_i\mu_i\,\dd n_i = \sum_i\mu_i\nu_i\,\dd\xi = \Delta_rG\,\dd\xi$. En égalant : $T\,\delta S_c = -\Delta_rG\,\dd\xi$.
\end{demo}

L'enthalpie libre joue donc, à $T$ et $P$ fixées, le rôle d'un \textbf{potentiel thermodynamique} : elle décroît au cours de toute évolution spontanée et elle est \emph{minimale à l'équilibre}. L'irréversibilité de la transformation est entièrement due à la réaction elle-même : $\delta S_c$ est non nulle tant que $\Delta_rG \neq 0$.

\section{Critère d'évolution spontanée}

\begin{resultat}{Sens d'évolution à $T$ et $P$ fixées}
Le deuxième principe impose $\Delta_rG\,\dd\xi \leq 0$ :
\begin{itemize}
  \item $\Delta_rG < 0 \;\Rightarrow\; \dd\xi > 0$ : évolution dans le \textbf{sens direct} ($\rightarrow$) ;
  \item $\Delta_rG > 0 \;\Rightarrow\; \dd\xi < 0$ : évolution dans le \textbf{sens indirect} ($\leftarrow$) ;
  \item $\Delta_rG = 0$ : \textbf{équilibre chimique}, aucune évolution.
\end{itemize}
\end{resultat}

\begin{pieges}
\begin{itemize}
  \item C'est le signe de $\Delta_rG$, \emph{pas} celui de $\Delta_rG^\circ$, qui donne le sens d'évolution. $\Delta_rG^\circ$ ne dépend que de $T$ ; $\Delta_rG$ dépend aussi de la composition actuelle.
  \item Le critère dit si une évolution est \emph{possible}, pas si elle est \emph{rapide} (exemple du diamant, chapitre~\ref{ch:potentiel}).
\end{itemize}
\end{pieges}

\section{Expression de l'enthalpie libre de réaction}

\begin{resultat}{Enthalpie libre de réaction et quotient de réaction}
\[ \boxed{\Delta_rG = \Delta_rG^\circ + RT\ln Q_r} \qquad\text{avec}\qquad Q_r = \prod_i a_i^{\,\nu_i} \]
\end{resultat}

\begin{demo}
On injecte $\mu_i = \mu_i^\circ + RT\ln a_i$ dans $\Delta_rG = \sum_i\nu_i\mu_i$ :
\[ \Delta_rG = \underbrace{\sum_i\nu_i\mu_i^\circ}_{\Delta_rG^\circ} + RT\sum_i\nu_i\ln a_i = \Delta_rG^\circ + RT\ln\Big(\prod_i a_i^{\,\nu_i}\Big) = \Delta_rG^\circ + RT\ln Q_r. \]
\end{demo}

\section{Constante d'équilibre et loi d'action des masses}

\begin{definition}[constante thermodynamique d'équilibre]
La \textbf{constante thermodynamique d'équilibre} d'une réaction à la température $T$ est
\[ \boxed{K^\circ(T) = \exp\left(-\frac{\Delta_rG^\circ(T)}{RT}\right)} \qquad\Longleftrightarrow\qquad \Delta_rG^\circ + RT\ln K^\circ = 0. \]
Elle est sans dimension et ne dépend que de $T$ (et de l'écriture de l'équation).
\end{definition}

\begin{resultat}{Loi d'action des masses (Guldberg et Waage) et critère en $Q_r$}
En combinant avec l'expression de $\Delta_rG$ :
\[ \boxed{\Delta_rG = RT\ln\frac{Q_r}{K^\circ}} \]
\begin{itemize}
  \item à l'équilibre, $\Delta_rG = 0$ donc $Q_{r,\text{éq}} = K^\circ(T)$ : c'est la \textbf{loi d'action des masses} ;
  \item hors équilibre, $Q_r < K^\circ \Leftrightarrow \Delta_rG < 0$ : sens direct ; $Q_r > K^\circ \Leftrightarrow \Delta_rG > 0$ : sens indirect.
\end{itemize}
\end{resultat}

\begin{figure}[H]
\centering
\begin{tikzpicture}[>=Stealth]
  \draw[->,line width=1pt] (0,0) -- (11,0) node[right] {$Q_r$};
  \draw (0,0.12) -- (0,-0.12) node[below] {$0$};
  \draw[line width=1.2pt,c3] (5.5,0.25) -- (5.5,-0.25) node[below,black] {$K^\circ$};
  \draw[->,line width=1.4pt,c4] (2.6,0.55) -- (4.6,0.55);
  \node[c4,font=\small,align=center] at (2.75,1.25) {$Q_r<K^\circ$ : $\Delta_rG<0$\\évolution dans le sens direct};
  \draw[->,line width=1.4pt,c2] (8.4,0.55) -- (6.4,0.55);
  \node[c2,font=\small,align=center] at (8.25,1.25) {$Q_r>K^\circ$ : $\Delta_rG>0$\\évolution dans le sens indirect};
  \node[c3,font=\small] at (5.5,-1.05) {équilibre : $Q_r = K^\circ$, $\Delta_rG = 0$};
\end{tikzpicture}
\caption{Le quotient de réaction évolue toujours vers la constante d'équilibre.}
\label{fig:axeQK}
\end{figure}

\begin{resultat}{Ce que dit le signe de $\Delta_rG^\circ$}
$K^\circ > 1 \Leftrightarrow \Delta_rG^\circ < 0$ : la réaction est dite \emph{thermodynamiquement favorisée} dans le sens direct (à partir de constituants dans l'état standard, $Q_r = 1 < K^\circ$). En pratique, si $K^\circ \gtrsim 10^4$, la réaction est souvent quasi totale ; si $K^\circ \lesssim 10^{-4}$, elle est quasi nulle. Ce ne sont que des repères : l'état final dépend aussi des quantités initiales.
\end{resultat}

\begin{resultat}{Les visages de $\Delta_rG^\circ$ -- à connaître}
\[ \Delta_rG^\circ = \sum_i\nu_i\mu_i^\circ = \sum_i\nu_i\,\Delta_fG^\circ_i = \Delta_rH^\circ - T\,\Delta_rS^\circ = -RT\ln K^\circ \]
et, pour une réaction rédox (chapitre~\ref{ch:redox}), $\Delta_rG^\circ = -\gamma F\,\Delta E^\circ$.
\end{resultat}

\begin{exemple}[ammoniac à $\SI{298}{K}$]
$\Delta_rG^\circ = \SI{-32,7}{kJ.mol^{-1}}$ (chapitre~\ref{ch:grandeurs}), donc
\[ K^\circ(\SI{298}{K}) = \exp\left(\frac{32{,}7\times10^3}{8{,}314\times298}\right) = \num{5,4e5}. \]
La synthèse est très favorisée à température ambiante\dots\ mais infiniment lente sans catalyseur ni chauffage (chapitre~\ref{ch:vanthoff}).
\end{exemple}

\section{Représentation graphique : la courbe $G(\xi)$}

À $T$ et $P$ fixées, $G$ ne dépend plus que de $\xi$. On peut la calculer avec $G = \sum_in_i\mu_i$ et le tableau d'avancement. La pente de la courbe est $\left(\partial G/\partial\xi\right)_{T,P} = \Delta_rG$.

\begin{exemple}[isomérisation en phase gazeuse]
\ce{A(g) = B(g)} sous $P = P^\circ$, avec $n_A = 1-\xi$, $n_B = \xi$ (en mol), donc $a_A = 1-\xi$ et $a_B = \xi$. En prenant $\mu_A^\circ = 0$ comme référence, $\mu_B^\circ = \Delta_rG^\circ$ :
\[ G(\xi) = \underbrace{\xi\,\Delta_rG^\circ}_{\text{réactifs et produits séparés}} + \underbrace{RT\big[(1-\xi)\ln(1-\xi) + \xi\ln\xi\big]}_{\text{terme de mélange}<0} \]
Avec $\Delta_rG^\circ = \SI{-2,0}{kJ.mol^{-1}}$ à $\SI{298}{K}$ : $K^\circ = 2{,}24$ et $\xi_{\text{éq}} = K^\circ/(1+K^\circ) = \SI{0,69}{mol}$.
\end{exemple}

\begin{figure}[H]
\centering
\begin{tikzpicture}
\begin{axis}[cours,width=0.76\linewidth,xlabel={$\xi$ (mol)},ylabel={$G - G(0)$ (kJ)},xmin=0,xmax=1.02,ymin=-3.3,ymax=0.4,
  legend style={at={(0.97,0.97)},anchor=north east}]
  \addplot[c1] table[x=xi,y=G]{data/gxi.dat}; \addlegendentry{$G(\xi)$}
  \addplot[black!50,dashed,line width=0.9pt] table[x=xi,y=Glin]{data/gxi.dat}; \addlegendentry{$\xi\,\Delta_rG^\circ$ (sans mélange)}
  \addplot[c4,line width=1.3pt] coordinates {(0.15,-1.499) (0.45,-2.729)};
  \addplot[c2,line width=1.3pt] coordinates {(0.80,-2.951) (1.00,-2.261)};
  \addplot[only marks,mark=*,mark size=1.8pt,c4] coordinates {(0.3,-2.114)};
  \addplot[only marks,mark=*,mark size=1.8pt,c2] coordinates {(0.9,-2.606)};
  \addplot[c3,dashed,line width=0.8pt] coordinates {(0.691,-3.3) (0.691,-2.915)};
  \addplot[only marks,mark=*,mark size=2.2pt,c3] coordinates {(0.691,-2.915)};
  \node[c4,font=\small,anchor=west] at (axis cs:0.02,-2.85) {$\Delta_rG<0$ : $\xi$ augmente};
  \node[c2,font=\small,anchor=east] at (axis cs:0.93,-2.18) {$\Delta_rG>0$ : $\xi$ diminue};
  \node[c3,font=\small,anchor=north east] at (axis cs:0.68,-3.0) {$\xi_{\text{éq}}$ : $\Delta_rG = 0$};
\end{axis}
\end{tikzpicture}
\caption{Courbe $G(\xi)$ : le système « descend » vers le minimum. La pente de la tangente est $\Delta_rG$.}
\label{fig:gxi}
\end{figure}

\noindent\textbf{Pourquoi le minimum est-il à l'intérieur ?} Le terme de mélange a une pente infinie en $\xi\to0$ et $\xi\to1$ ($RT\ln\xi\to-\infty$) : dès qu'un constituant d'une phase mélange (gaz, soluté) a une quantité nulle, $\Delta_rG$ est infini et pousse le système à le former. Un équilibre entre gaz ou solutés n'est donc jamais rigoureusement total. Il en va autrement si les seuls constituants qui disparaissent sont des solides ou liquides \emph{purs} : leur activité reste égale à 1 jusqu'à leur disparition, et la réaction peut s'arrêter par épuisement avant d'atteindre l'équilibre.

\section{État final : équilibre ou transformation totale}

\begin{definition}[équilibre chimique, transformation totale, rupture d'équilibre]
\begin{itemize}
  \item L'état final est un \textbf{équilibre chimique} si tous les constituants de l'équation sont présents : alors $\Delta_rG = 0$ et $Q_r = K^\circ$.
  \item La transformation est \textbf{totale} si un réactif est entièrement consommé : $\xi_F = \xi_{\max}$. Le système s'arrête alors \emph{hors équilibre}, avec $\Delta_rG < 0$ ($Q_r < K^\circ$), faute de réactif.
  \item On parle de \textbf{rupture d'équilibre} quand l'espèce épuisée est une phase condensée pure (solide, liquide), ou un précipité qui se dissout complètement.
\end{itemize}
\end{definition}

\begin{methode}[déterminer l'état final]
\begin{enumerate}
  \item Tableau d'avancement ; calculer $\xi_{\max}$ (et $\xi_{\min}$ si le sens n'est pas connu).
  \item Calculer $Q_{r,0}$ et le comparer à $K^\circ$ pour connaître le sens d'évolution.
  \item \emph{Supposer} l'équilibre atteint : résoudre $Q_r(\xi_{\text{éq}}) = K^\circ$.
  \item \emph{Vérifier} : si $\xi_{\text{éq}}$ est accessible ($\xi_{\text{éq}}\leq\xi_{\max}$), l'hypothèse est validée. Sinon, la transformation est totale : $\xi_F = \xi_{\max}$, et l'on vérifie $Q_r(\xi_{\max}) < K^\circ$.
  \item Si $K^\circ$ est très grande, on peut supposer d'emblée la réaction quasi totale, puis calculer la très faible quantité restante du réactif limitant à l'aide de $K^\circ$.
\end{enumerate}
\end{methode}

\begin{exemple}[décomposition du calcaire en enceinte fermée]
\ce{CaCO3(s) = CaO(s) + CO2(g)}, $K^\circ = P_{\ce{CO2}}/P^\circ$ (les solides ont une activité 1). Données : $\Delta_rH^\circ = \SI{179,2}{kJ.mol^{-1}}$, $\Delta_rS^\circ = \SI{160,2}{J.K^{-1}.mol^{-1}}$.

\noindent À $T = \SI{1100}{K}$ : $\Delta_rG^\circ = 179{,}2 - 1100\times0{,}1602 = \SI{3,0}{kJ.mol^{-1}}$ et $K^\circ = \mathrm{e}^{-3{,}0\times10^3/(8{,}314\times1100)} = 0{,}722$.

\noindent On introduit $n_0$ de \ce{CaCO3} dans une enceinte vide de $V = \SI{10}{L}$. Si l'équilibre est atteint, $P_{\ce{CO2}} = \SI{0,722}{bar}$, ce qui demande
\[ \xi_{\text{éq}} = \frac{P_{\ce{CO2}}V}{RT} = \frac{0{,}722\times10^5\times10^{-2}}{8{,}314\times1100} = \SI{78,9}{mmol}. \]
\begin{itemize}
  \item $n_0 = \SI{0,100}{mol} > \xi_{\text{éq}}$ : équilibre, il reste $\SI{21}{mmol}$ de \ce{CaCO3}, $P = \SI{0,722}{bar}$.
  \item $n_0 = \SI{0,050}{mol} < \xi_{\text{éq}}$ : \textbf{rupture d'équilibre}. Tout le calcaire se décompose ; $P = \dfrac{n_0RT}{V} = \SI{0,457}{bar} < K^\circ P^\circ$, et $\Delta_rG = RT\ln(Q_r/K^\circ) = \SI{-4,2}{kJ.mol^{-1}} < 0$ : la réaction « voudrait » continuer mais il n'y a plus de réactif.
\end{itemize}
\end{exemple}

\begin{figure}[H]
\centering
\begin{tikzpicture}
\begin{axis}[cours,width=0.72\linewidth,height=0.4\linewidth,xlabel={$n_0(\ce{CaCO3})$ (mmol)},ylabel={$P_{\ce{CO2}}$ final (bar)},
  xmin=0,xmax=150,ymin=0,ymax=0.9]
  \addplot[c1] table[x=n0_mmol,y=P_bar]{data/caco3.dat};
  \addplot[c3,dashed,line width=0.8pt] coordinates {(78.9,0) (78.9,0.722)};
  \node[c3,font=\small,anchor=north west] at (axis cs:80,0.18) {$\xi_{\text{éq}} = \SI{78,9}{mmol}$};
  \node[font=\small,align=center,c2] at (axis cs:36,0.68) {transformation\\totale};
  \node[font=\small,align=center,c4] at (axis cs:117,0.55) {équilibre\\$P = K^\circ P^\circ$};
\end{axis}
\end{tikzpicture}
\caption{Pression finale en \ce{CO2} à $\SI{1100}{K}$ dans $\SI{10}{L}$ en fonction de la quantité initiale de calcaire.}
\label{fig:caco3}
\end{figure}

% ================================================================== Chapitre 6
\chapter{Influence des paramètres et optimisation}
\label{ch:vanthoff}

Une fois l'équilibre établi, on peut le \emph{déplacer} en jouant sur les conditions opératoires. Deux leviers existent, et toute la discussion s'organise autour d'eux : modifier $K^\circ$ (seule la température le permet) ou modifier $Q_r$ (pression, composition).

\section{Influence de la température : relation de Van~'t~Hoff}

\begin{resultat}{Relation de Van~'t~Hoff}
\[ \boxed{\frac{\dd\ln K^\circ}{\dd T} = \frac{\Delta_rH^\circ}{RT^2}} \]
\begin{itemize}
  \item réaction \textbf{endothermique} ($\Delta_rH^\circ > 0$) : $K^\circ$ \emph{augmente} avec $T$ ;
  \item réaction \textbf{exothermique} ($\Delta_rH^\circ < 0$) : $K^\circ$ \emph{diminue} quand $T$ augmente ;
  \item réaction athermique : $K^\circ$ quasi indépendante de $T$.
\end{itemize}
\end{resultat}

\begin{demo}[Démonstration dans l'approximation d'Ellingham]
$\ln K^\circ = -\dfrac{\Delta_rG^\circ}{RT} = -\dfrac{\Delta_rH^\circ}{RT} + \dfrac{\Delta_rS^\circ}{R}$. Avec $\Delta_rH^\circ$ et $\Delta_rS^\circ$ constantes, la dérivée par rapport à $T$ vaut $\dfrac{\Delta_rH^\circ}{RT^2}$.
\end{demo}

\begin{horsprog}[validité générale]
La relation de Van~'t~Hoff reste vraie même si $\Delta_rH^\circ$ dépend de $T$. On le montre avec la relation de Gibbs--Helmholtz $\dfrac{\dd}{\dd T}\left(\dfrac{\Delta_rG^\circ}{T}\right) = -\dfrac{\Delta_rH^\circ}{T^2}$, conséquence de $\left(\partial G/\partial T\right)_P = -S$.
\end{horsprog}

\begin{resultat}{Forme intégrée ($\Delta_rH^\circ$ constante)}
\[ \ln\frac{K^\circ(T_2)}{K^\circ(T_1)} = -\frac{\Delta_rH^\circ}{R}\left(\frac{1}{T_2}-\frac{1}{T_1}\right) \]
$\ln K^\circ$ est une fonction affine de $1/T$, de pente $-\Delta_rH^\circ/R$.
\end{resultat}

\begin{figure}[H]
\centering
\begin{tikzpicture}
\begin{axis}[cours,width=0.74\linewidth,height=0.42\linewidth,xlabel={$1000/T$ ($\mathrm{K^{-1}}$)},ylabel={$\ln K^\circ$},
  xmin=1.0,xmax=2.6,ymin=-12,ymax=3,clip=false]
  \addplot[c1,mark=*,mark size=1.6pt] table[x=invT_milli,y=lnK]{data/nh3_vh.dat};
  \draw[c3,line width=0.8pt] (axis cs:1.6,-6.16) -- (axis cs:2.0,-6.16) -- (axis cs:2.0,-1.74);
  \node[c3,font=\small,anchor=west,align=left] at (axis cs:2.04,-5.6) {pente $= -\Delta_rH^\circ/R$\\$= +\SI{1,10e4}{K}$};
  \node[font=\small,anchor=south east,align=right] at (axis cs:2.58,-11.6) {\ce{N2 + 3 H2 = 2 NH3}, exothermique :\\$K^\circ$ diminue quand $T$ augmente};
\end{axis}
\end{tikzpicture}
\caption{Droite de Van~'t~Hoff pour la synthèse de l'ammoniac (approximation d'Ellingham), de $\SI{400}{K}$ à $\SI{900}{K}$.}
\label{fig:vh}
\end{figure}

\begin{methode}[calculer $K^\circ$ à une température quelconque]
\begin{itemize}
  \item \emph{Voie 1} : calculer $\Delta_rG^\circ(T) = \Delta_rH^\circ - T\Delta_rS^\circ$ avec les tables, puis $K^\circ(T) = \exp(-\Delta_rG^\circ/RT)$.
  \item \emph{Voie 2} : partir d'une valeur connue $K^\circ(T_1)$ et utiliser la forme intégrée de Van~'t~Hoff.
\end{itemize}
Les deux voies donnent le même résultat dans l'approximation d'Ellingham.
\end{methode}

\begin{exemple}[ammoniac à $\SI{450}{\celsius}$]
Voie 1 : $\Delta_rG^\circ(\SI{723}{K}) = -91{,}8 + 723\times0{,}1981 = \SI{51,4}{kJ.mol^{-1}}$, d'où $K^\circ = \mathrm{e}^{-51{,}4\times10^3/(8{,}314\times723)} = \num{1,9e-4}$.

\noindent Voie 2 : $K^\circ(723) = \num{5,4e5}\times\exp\left[\dfrac{91{,}8\times10^3}{8{,}314}\left(\dfrac{1}{723}-\dfrac{1}{298}\right)\right] = \num{1,9e-4}$.

\noindent En passant de $\SI{25}{\celsius}$ à $\SI{450}{\celsius}$, $K^\circ$ est divisée par environ $3\times10^9$.
\end{exemple}

\begin{methode}[mesurer une enthalpie standard de réaction]
On mesure $K^\circ$ à plusieurs températures (par titrage, conductimétrie, mesure de pression\dots), on trace $\ln K^\circ$ en fonction de $1/T$ et on fait une régression linéaire : la pente vaut $-\Delta_rH^\circ/R$ et l'ordonnée à l'origine $\Delta_rS^\circ/R$. C'est aussi une manière de vérifier l'approximation d'Ellingham (les points sont-ils alignés ?).
\end{methode}

\section{Optimisation par modification de $K^\circ$ : choix de la température}

Seule la température modifie $K^\circ$. Pour \emph{augmenter le rendement} à l'équilibre :
\begin{itemize}
  \item réaction exothermique : travailler à \emph{basse} température ;
  \item réaction endothermique : travailler à \emph{haute} température.
\end{itemize}
Mais la vitesse de réaction augmente toujours avec $T$ (loi d'Arrhenius). Pour une réaction exothermique, il y a donc un \textbf{compromis} entre thermodynamique (basse $T$) et cinétique (haute $T$), que l'on améliore avec un catalyseur (qui accélère la réaction sans modifier $K^\circ$).

\section{Optimisation par modification de $Q_r$}

On part d'un système à l'équilibre ($Q_r = K^\circ$). On le perturbe à $T$ fixée, \emph{instantanément}, puis on regarde comment $Q_r$ a varié : le système évolue ensuite de façon à ramener $Q_r$ vers $K^\circ$.

\begin{resultat}{Quotient de réaction d'un mélange gazeux}
Pour des gaz parfaits, $a_i = \dfrac{P_i}{P^\circ} = \dfrac{n_i}{n_{\text{tot,gaz}}}\,\dfrac{P}{P^\circ}$, d'où, en notant $\Delta\nu_{\text{gaz}} = \sum_{\text{gaz}}\nu_i$ :
\[ Q_r = \Big(\prod_{\text{gaz}}n_i^{\,\nu_i}\Big)\left(\frac{P}{n_{\text{tot,gaz}}\,P^\circ}\right)^{\Delta\nu_{\text{gaz}}}\times\big(\text{activités des autres constituants}\big). \]
\end{resultat}

\begin{methode}[prévoir l'effet d'une perturbation]
\begin{enumerate}
  \item Exprimer $Q_r$ en fonction de la grandeur modifiée ($P$, $n_i$, $n_{\text{tot}}$\dots), les autres restant fixes.
  \item Dire si la perturbation fait augmenter ou diminuer $Q_r$ à partir de $Q_r = K^\circ$.
  \item Conclure : $Q_r < K^\circ$ entraîne une évolution dans le sens direct ; $Q_r > K^\circ$ dans le sens indirect.
\end{enumerate}
\end{methode}

\begin{resultat}{Effet des principaux paramètres (gaz parfaits, $T$ fixée)}
\renewcommand{\arraystretch}{1.35}
\begin{tabularx}{\linewidth}{>{\raggedright\arraybackslash}p{4.4cm} X}
\toprule
\textbf{Perturbation} & \textbf{Effet sur $Q_r$ et conséquence} \\ \midrule
augmentation de la pression totale $P$ & $Q_r\propto P^{\Delta\nu_{\text{gaz}}}$ : si $\Delta\nu_{\text{gaz}}<0$, $Q_r$ diminue $\Rightarrow$ sens direct (vers \emph{moins} de gaz) ; si $\Delta\nu_{\text{gaz}} = 0$, aucun effet \\
ajout d'un gaz inerte à $P$ fixée & $n_{\text{tot,gaz}}$ augmente : $Q_r\propto n_{\text{tot,gaz}}^{-\Delta\nu_{\text{gaz}}}$. Si $\Delta\nu_{\text{gaz}}<0$, $Q_r$ augmente $\Rightarrow$ sens indirect (même effet qu'une baisse de $P$) \\
ajout d'un gaz inerte à $V$ fixé & les pressions partielles $P_i = n_iRT/V$ sont inchangées : aucun effet \\
élimination d'un produit (distillation, liquéfaction, précipitation) & $Q_r$ diminue $\Rightarrow$ sens direct \\
ajout d'un réactif (en solution, ou gaz à $V$ fixé) & $Q_r$ diminue $\Rightarrow$ sens direct \\
ajout ou retrait d'un solide ou liquide pur (sans le faire disparaître) & activité égale à 1, $Q_r$ inchangé : aucun effet \\
\bottomrule
\end{tabularx}
\end{resultat}

\begin{pieges}
\begin{itemize}
  \item Raisonner sur $Q_r$, pas sur des intuitions : l'ajout d'un gaz inerte n'a pas le même effet à $P$ fixée et à $V$ fixé.
  \item Un catalyseur ne modifie ni $K^\circ$ ni $Q_r$ : il ne change pas l'état d'équilibre, il permet seulement de l'atteindre plus vite.
\end{itemize}
\end{pieges}

\begin{horsprog}[principe de modération de Le Chatelier]
Les résultats du tableau se résument qualitativement ainsi : « le système réagit à une perturbation en s'opposant à celle-ci » (une hausse de $P$ favorise le sens qui consomme du gaz, une hausse de $T$ favorise le sens endothermique). C'est un moyen mnémotechnique ; aux concours, on attend la justification par $Q_r$ et $K^\circ$.
\end{horsprog}

\section{Application : la synthèse industrielle de l'ammoniac}

Le procédé Haber--Bosch réalise \ce{N2(g) + 3 H2(g) = 2 NH3(g)}, avec $\Delta_rH^\circ = \SI{-91,8}{kJ.mol^{-1}}$ (exothermique) et $\Delta\nu_{\text{gaz}} = -2$.

\begin{itemize}
  \item \textbf{Température} : la thermodynamique demande une basse température, mais la réaction est alors trop lente. Compromis : 400 à $\SI{500}{\celsius}$ avec un catalyseur à base de fer.
  \item \textbf{Pression} : $\Delta\nu_{\text{gaz}}<0$, donc une haute pression favorise la synthèse. On travaille vers 150 à $\SI{300}{bar}$ (au-delà, le coût des installations devient prohibitif).
  \item \textbf{Composition} : on élimine \ce{NH3} par liquéfaction à la sortie du réacteur, et l'on recycle \ce{N2} et \ce{H2} non transformés. On évite l'accumulation de gaz inertes (argon de l'air), qui défavorise la synthèse à $P$ fixée.
\end{itemize}

\begin{exemple}[effet de la pression à $\SI{723}{K}$]
Mélange initial stœchiométrique (\SI{1}{mol} de \ce{N2}, \SI{3}{mol} de \ce{H2}) ; on résout $Q_r(\xi) = K^\circ = \num{1,9e-4}$ avec $n_{\text{tot}} = 4 - 2\xi$ :
\[ \frac{(2\xi)^2(4-2\xi)^2}{(1-\xi)(3-3\xi)^3}\left(\frac{P^\circ}{P}\right)^2 = K^\circ \]
\begin{center}
\renewcommand{\arraystretch}{1.25}
\begin{tabular}{lcccc}
\toprule
$P$ (bar) & 1 & 50 & 200 & 300 \\ \midrule
$\xi_{\text{éq}}$ (mol) & 0,009 & 0,27 & 0,53 & 0,60 \\
fraction molaire de \ce{NH3} & 0,004 & 0,16 & 0,36 & 0,43 \\
\bottomrule
\end{tabular}
\end{center}
Sous $\SI{1}{bar}$ à $\SI{450}{\celsius}$, la synthèse est quasi nulle ; à $\SI{200}{bar}$, plus du tiers du mélange est de l'ammoniac. (Valeurs calculées numériquement par résolution de l'équation ci-dessus.)
\end{exemple}

% ================================================================== Chapitre 7
\chapter{Thermodynamique de l'oxydoréduction}
\label{ch:redox}

Les potentiels d'électrode de la première année et l'enthalpie libre de réaction décrivent la même chose : la tendance d'un système rédox à évoluer. Ce chapitre fait le lien, ce qui permet de calculer des potentiels standard à partir de tables thermodynamiques, et inversement.

\section{Enthalpie libre d'une réaction d'oxydoréduction}

On considère la réaction entre le couple 1 ($\mathrm{Ox}_1/\mathrm{Red}_1$) et le couple 2 ($\mathrm{Ox}_2/\mathrm{Red}_2$), écrite pour échanger $\gamma$ électrons :
\[ (\mathcal{R})\qquad a_1\,\mathrm{Ox}_1 + b_2\,\mathrm{Red}_2 \;=\; b_1\,\mathrm{Red}_1 + a_2\,\mathrm{Ox}_2 \]

\begin{resultat}{Relation entre enthalpie libre de réaction et potentiels}
\[ \boxed{\Delta_rG = -\gamma F\,(E_1 - E_2)} \qquad\qquad \boxed{\Delta_rG^\circ = -\gamma F\,(E_1^\circ - E_2^\circ)} \]
$E_1$ et $E_2$ sont les potentiels de Nernst des deux couples calculés avec la composition \emph{actuelle} du système ; $F$ est la constante de Faraday.
\end{resultat}

\noindent\textbf{D'où vient le facteur $-\gamma F$ ?} Il est établi au chapitre~\ref{ch:piles} : lorsque la réaction se déroule dans une pile, l'avancement $\dd\xi$ fait circuler une charge $\gamma F\,\dd\xi$ sous la tension $E_1 - E_2$, et $\Delta_rG\,\dd\xi$ est égale au travail électrique reçu. On vérifie ici que la relation est \emph{cohérente} avec la loi de Nernst.

\begin{demo}[Cohérence avec la formule de Nernst]
En écrivant Nernst pour chaque couple (demi-équations multipliées pour faire intervenir $\gamma$ électrons) et en soustrayant :
\[ E_1 - E_2 = (E_1^\circ - E_2^\circ) + \frac{RT}{\gamma F}\ln\frac{a_{\mathrm{Ox}_1}^{a_1}\,a_{\mathrm{Red}_2}^{b_2}}{a_{\mathrm{Red}_1}^{b_1}\,a_{\mathrm{Ox}_2}^{a_2}} = (E_1^\circ - E_2^\circ) - \frac{RT}{\gamma F}\ln Q_r. \]
En multipliant par $-\gamma F$ : $-\gamma F(E_1-E_2) = -\gamma F(E_1^\circ-E_2^\circ) + RT\ln Q_r$. C'est exactement la structure $\Delta_rG = \Delta_rG^\circ + RT\ln Q_r$ : les deux relations encadrées sont compatibles entre elles et avec le chapitre~\ref{ch:evolution}.
\end{demo}

\begin{resultat}{Conséquences}
\begin{itemize}
  \item \textbf{Sens d'évolution} : $\Delta_rG < 0 \Leftrightarrow E_1 > E_2$. L'oxydant du couple de plus haut potentiel est réduit : c'est la justification thermodynamique de la règle du gamma.
  \item \textbf{Équilibre} : $\Delta_rG = 0 \Leftrightarrow E_1 = E_2$ : la solution n'a plus qu'un seul potentiel.
  \item \textbf{Constante d'équilibre} : $-RT\ln K^\circ = -\gamma F(E_1^\circ - E_2^\circ)$, d'où
  \[ \ln K^\circ = \frac{\gamma F}{RT}(E_1^\circ - E_2^\circ) \qquad\Longrightarrow\qquad \log K^\circ \underset{\SI{298}{K}}{=} \frac{\gamma(E_1^\circ - E_2^\circ)}{0{,}06}. \]
  \item Écriture compacte : $\Delta_rG = RT\ln\dfrac{Q_r}{K^\circ} = -\gamma F\,\Delta E$, avec $\Delta E = E_1 - E_2$.
\end{itemize}
\end{resultat}

\begin{ordres}
À $\SI{298}{K}$, un écart $\Delta E^\circ = \SI{0,1}{V}$ pour $\gamma = 1$ correspond à $\Delta_rG^\circ = \SI{-9,6}{kJ.mol^{-1}}$ et à $K^\circ\approx 10^{1{,}7}\approx 50$. Pour $\gamma = 2$ et $\Delta E^\circ = \SI{1,10}{V}$ (pile Daniell) : $\Delta_rG^\circ = \SI{-212}{kJ.mol^{-1}}$ et $K^\circ \approx 10^{37}$.
\end{ordres}

\section{Enthalpie libre standard d'une demi-équation}

\begin{definition}[enthalpie libre standard de demi-réaction]
À la demi-équation $\alpha\,\mathrm{Ox} + \gamma\,e^- + \dots = \beta\,\mathrm{Red} + \dots$, on associe par convention
\[ \boxed{\Delta_{1/2}G^\circ = -\gamma F\,E^\circ(\mathrm{Ox/Red})} \]
\end{definition}

Cette convention revient à mesurer chaque couple par rapport à l'ESH : $\Delta_{1/2}G^\circ$ est l'enthalpie libre standard de la réaction complète entre $\mathrm{Ox}$ et \ce{H2} (couple \ce{H+/H2}, de potentiel standard nul). Son intérêt :
\begin{itemize}
  \item pour toute réaction rédox, $\Delta_rG^\circ = \Delta_{1/2}G_1^\circ - \Delta_{1/2}G_2^\circ$ : les électrons « s'éliminent » comme dans l'équation ;
  \item les $\Delta_{1/2}G^\circ$ s'additionnent comme des grandeurs de réaction ordinaires, et l'on peut les combiner avec les $\Delta_rG^\circ$ d'autres réactions (dissolution, acido-basicité, formation).
\end{itemize}

\section{Calcul d'un potentiel standard inconnu}

\begin{methode}[potentiel standard d'un couple à partir de données]
\begin{enumerate}
  \item Écrire la demi-équation cherchée, avec ses $\gamma$ électrons.
  \item L'obtenir comme \emph{combinaison linéaire} de demi-équations connues et de réactions de constante connue ($K_s$, $K_a$, $K_e$, $\Delta_fG^\circ$\dots).
  \item Additionner les enthalpies libres standard : $-\gamma F E^\circ_{\text{cherché}} = \sum(\text{contributions})$, avec $\Delta_{1/2}G^\circ = -\gamma_kFE^\circ_k$ pour une demi-équation et $\Delta_rG^\circ = -RT\ln K^\circ$ pour une réaction.
  \item Diviser par $-\gamma F$.
\end{enumerate}
On n'additionne \textbf{jamais} les potentiels : ce sont les enthalpies libres qui s'additionnent.
\end{methode}

\subsection{Même nombre d'oxydation dans les deux couples}

\begin{exemple}[électrode au calomel : $E^\circ(\ce{Hg2Cl2}/\ce{Hg})$]
Données : \ce{Hg2^{2+} + 2e^- = 2 Hg(l)}, $E_1^\circ = \SI{0,797}{V}$ ; \ce{Hg2Cl2(s) = Hg2^{2+} + 2 Cl^-}, $\mathrm{p}K_s = 17{,}9$. Dans les deux couples, le mercure passe du n.o. $+\mathrm{I}$ à $0$ : seule l'espèce oxydée diffère (ion libre ou précipité).

\medskip
\noindent\emph{Méthode 1 -- égalité des potentiels.} Une solution au contact de \ce{Hg}, \ce{Hg2Cl2} et \ce{Cl-} n'a qu'un seul potentiel, donné par l'un ou l'autre couple :
\[ E = E_1^\circ + \frac{0{,}059}{2}\log\frac{[\ce{Hg2^{2+}}]}{c^\circ} \qquad\text{avec}\qquad \frac{[\ce{Hg2^{2+}}]}{c^\circ} = \frac{K_s}{([\ce{Cl-}]/c^\circ)^2} \]
\[ \Longrightarrow\quad E = \underbrace{E_1^\circ - \frac{0{,}059}{2}\,\mathrm{p}K_s}_{E^\circ(\ce{Hg2Cl2/Hg})} - 0{,}059\log\frac{[\ce{Cl-}]}{c^\circ}. \]
\emph{Méthode 2 -- combinaison des enthalpies libres.}
\[ \renewcommand{\arraystretch}{1.3}\begin{array}{cl@{\qquad}l}
 & \ce{Hg2Cl2(s) -> Hg2^{2+} + 2Cl-} & \Delta_rG^\circ = -RT\ln K_s\\
 + & \ce{Hg2^{2+} + 2e^- -> 2Hg(l)} & \Delta_{1/2}G^\circ = -2FE_1^\circ\\ \hline
 = & \ce{Hg2Cl2(s) + 2e^- -> 2Hg(l) + 2Cl-} & -2FE^\circ = -RT\ln K_s - 2FE_1^\circ
\end{array} \]
d'où $E^\circ = E_1^\circ + \dfrac{RT}{2F}\ln K_s = 0{,}797 - 0{,}0296\times17{,}9 = \SI{0,27}{V}$.
\end{exemple}

\begin{pieges}
Ici le terme $\frac{0{,}059}{2}\mathrm{p}K_s$ vaut $\SI{0,53}{V}$ : arrondir $0{,}059$ à $0{,}06$ introduit une erreur de $\SI{0,01}{V}$. Quand un coefficient est multiplié par un grand $\mathrm{p}K$, garder $0{,}059$.
\end{pieges}

\subsection{Nombres d'oxydation différents}

\begin{exemple}[$E^\circ(\ce{Fe^{3+}}/\ce{Fe})$]
Données : \ce{Fe^{3+} + e^- = Fe^{2+}}, $E_1^\circ = \SI{0,77}{V}$ ; \ce{Fe^{2+} + 2e^- = Fe(s)}, $E_2^\circ = \SI{-0,44}{V}$. La demi-équation cherchée \ce{Fe^{3+} + 3e^- = Fe(s)} est la somme des deux :
\[ -3F\,E_3^\circ = -1\cdot F\,E_1^\circ - 2F\,E_2^\circ \quad\Longrightarrow\quad E_3^\circ = \frac{E_1^\circ + 2E_2^\circ}{3} = \frac{0{,}77 - 0{,}88}{3} = \SI{-0,04}{V}. \]
\end{exemple}

\begin{figure}[H]
\centering
\begin{tikzpicture}[>=Stealth]
  \node[font=\large] (A) at (0,0) {\ce{Fe^{3+}}};
  \node[font=\large] (B) at (4,0) {\ce{Fe^{2+}}};
  \node[font=\large] (C) at (8,0) {\ce{Fe(s)}};
  \draw[->,line width=1.1pt,navy] (A) -- node[below,font=\small] {$\gamma=1$} node[above,font=\small] {$\SI{0,77}{V}$} (B);
  \draw[->,line width=1.1pt,navy] (B) -- node[below,font=\small] {$\gamma=2$} node[above,font=\small] {$\SI{-0,44}{V}$} (C);
  \draw[->,line width=1.1pt,c3] (A.north) to[out=35,in=145] node[above,font=\small] {$\gamma = 3$ : $E^\circ = \dfrac{1\times0{,}77 + 2\times(-0{,}44)}{3} = \SI{-0,04}{V}$} (C.north);
\end{tikzpicture}
\caption{Diagramme de potentiels du fer : les potentiels se combinent pondérés par le nombre d'électrons.}
\label{fig:latimer}
\end{figure}

\subsection{À partir des enthalpies libres de formation}

\begin{exemple}[$E^\circ(\ce{O2}/\ce{H2O})$]
Demi-équation : \ce{O2(g) + 4 H+ + 4 e^- = 2 H2O(l)}. Associée à l'ESH, elle correspond à la réaction \ce{O2(g) + 2 H2(g) -> 2 H2O(l)} échangeant $\gamma = 4$ électrons, avec $E^\circ(\ce{H+/H2}) = 0$ :
\[ \Delta_rG^\circ = 2\,\Delta_fG^\circ(\ce{H2O(l)}) = 2\times(-237{,}1) = \SI{-474,2}{kJ.mol^{-1}} = -4F\,(E^\circ - 0) \]
\[ \Longrightarrow\quad E^\circ(\ce{O2/H2O}) = \frac{474{,}2\times10^3}{4\times96485} = \SI{1,23}{V}. \]
\end{exemple}

% ================================================================== Chapitre 8
\chapter{Piles électrochimiques}
\label{ch:piles}

Une pile convertit l'énergie libérée par une réaction rédox spontanée en énergie électrique. L'enthalpie libre fixe à la fois la tension qu'elle peut délivrer et l'énergie maximale qu'elle peut fournir.

\section{Description d'une pile}

\begin{definition}[pile électrochimique et vocabulaire]
\begin{itemize}
  \item Une \textbf{pile} est un générateur électrique dans lequel une réaction d'oxydoréduction spontanée a lieu \emph{sans contact direct} entre les réactifs : l'oxydation et la réduction se déroulent dans deux \textbf{demi-piles} séparées, reliées par une \textbf{jonction électrolytique} (pont salin, paroi poreuse) qui laisse passer les ions sans mélanger les solutions. Les électrons passent par le circuit extérieur.
  \item L'\textbf{anode} est l'électrode où a lieu l'\emph{oxydation} ; la \textbf{cathode} celle où a lieu la \emph{réduction}. Dans une pile, l'anode est le pôle $\ominus$ et la cathode le pôle $\oplus$.
  \item La \textbf{tension à vide} $e$ (anciennement « force électromotrice ») est la tension $V_\oplus - V_\ominus$ mesurée en circuit ouvert ($i = 0$). Elle vaut $e = E_\oplus - E_\ominus = E_{\text{cathode}} - E_{\text{anode}} > 0$.
\end{itemize}
\end{definition}

\begin{figure}[H]
\centering
\begin{tikzpicture}[>=Stealth,scale=0.95]
  % béchers
  \draw[line width=1pt] (0,0) -- (0,3) (0,0) -- (3,0) -- (3,3);
  \draw[line width=1pt] (6,0) -- (6,3) (6,0) -- (9,0) -- (9,3);
  \fill[c1!10] (0.03,0.03) rectangle (2.97,2.3);
  \fill[c1!25] (6.03,0.03) rectangle (8.97,2.3);
  % électrodes
  \fill[black!45] (1.2,0.6) rectangle (1.6,3.7); \node[font=\small,anchor=east] at (1.15,3.4) {Zn};
  \fill[c2!80!black] (7.4,0.6) rectangle (7.8,3.7); \node[font=\small,anchor=west] at (7.85,3.4) {Cu};
  \node[font=\small,align=center] at (0.75,1.0) {\ce{Zn^{2+}}\\\ce{SO4^{2-}}};
  \node[font=\small,align=center] at (8.35,1.0) {\ce{Cu^{2+}}\\\ce{SO4^{2-}}};
  % pont salin
  \draw[line width=6pt,black!20] (2.3,1.4) -- (2.3,2.9) -- (6.7,2.9) -- (6.7,1.4);
  \node[font=\scriptsize] at (4.5,3.2) {pont salin (\ce{K+}, \ce{NO3-})};
  \draw[->,c3,line width=0.9pt] (4.85,2.45) -- node[below,font=\scriptsize] {\ce{K+}} (5.95,2.45);
  \draw[->,c1,line width=0.9pt] (4.15,2.45) -- node[below,font=\scriptsize] {\ce{NO3-}} (3.05,2.45);
  % circuit extérieur
  \draw[line width=0.9pt] (1.4,3.7) -- (1.4,5.2) -- (3.7,5.2);
  \draw[line width=0.9pt] (5.3,5.2) -- (7.6,5.2) -- (7.6,3.7);
  \draw[line width=0.9pt,fill=white] (3.7,4.95) rectangle (5.3,5.45); \node[font=\small] at (4.5,5.2) {$R$};
  \draw[->,navy,line width=1.1pt] (2.0,5.55) -- (3.4,5.55) node[midway,above,font=\small] {$e^-$};
  \draw[->,c3,line width=1.1pt] (7.0,4.85) -- (5.6,4.85) node[midway,below,font=\small] {$i$};
  % pôles
  \node[font=\large] at (0.9,4.5) {$\ominus$}; \node[font=\large] at (8.1,4.5) {$\oplus$};
  % légendes
  \node[font=\small,align=center] at (1.5,-0.8) {\textbf{anode} : oxydation\\\ce{Zn -> Zn^{2+} + 2e^-}};
  \node[font=\small,align=center] at (7.5,-0.8) {\textbf{cathode} : réduction\\\ce{Cu^{2+} + 2e^- -> Cu}};
\end{tikzpicture}
\caption{Pile Daniell en fonctionnement. Les électrons circulent de l'anode vers la cathode dans le circuit extérieur ; dans le pont salin, les cations migrent vers la cathode et les anions vers l'anode pour maintenir l'électroneutralité.}
\label{fig:daniell}
\end{figure}

\noindent Notation conventionnelle : $\ominus\ \ce{Zn(s)} \,|\, \ce{Zn^{2+}},\ce{SO4^{2-}} \,\|\, \ce{Cu^{2+}},\ce{SO4^{2-}} \,|\, \ce{Cu(s)}\ \oplus$. Réaction de fonctionnement : \ce{Zn(s) + Cu^{2+} -> Zn^{2+} + Cu(s)}, avec $\gamma = 2$.

\section{Bilan énergétique : enthalpie libre et travail électrique}

\begin{resultat}{Inégalité entre variation d'enthalpie libre et travail électrique}
Pour un système évoluant de façon monotherme et monobare ($T = T_{\text{ext}}$, $P = P_{\text{ext}}$ aux états initial et final), recevant un travail électrique $W_{\text{élec}}$ en plus du travail des forces de pression :
\[ \boxed{\Delta G \leq W_{\text{élec}}} \qquad\text{(égalité si la transformation est réversible).} \]
Pour une pile, $W_{\text{élec}} < 0$ : l'énergie électrique \emph{fournie} au circuit vérifie
\[ W_{\text{fournie}} = -W_{\text{élec}} \leq -\Delta G, \qquad\text{donc}\qquad W_{\max} = -\Delta G. \]
\end{resultat}

\begin{demo}
Sur une évolution élémentaire, avec $\delta W = -P\,\dd V + \delta W_{\text{élec}}$ :
\begin{itemize}
  \item premier principe : $\dd U = -P\,\dd V + \delta W_{\text{élec}} + \delta Q$, donc $\dd H = \dd U + P\,\dd V = \delta W_{\text{élec}} + \delta Q$ (à $P$ fixée) ;
  \item deuxième principe : $\dd S = \dfrac{\delta Q}{T} + \delta S_c$ avec $\delta S_c \geq 0$ ;
  \item à $T$ fixée : $\dd G = \dd H - T\,\dd S = \delta W_{\text{élec}} - T\,\delta S_c \leq \delta W_{\text{élec}}$.
\end{itemize}
En intégrant : $\Delta G \leq W_{\text{élec}}$, l'égalité correspondant à $\delta S_c = 0$.
\end{demo}

\noindent L'enthalpie libre porte bien son nom : $-\Delta G$ est la part de l'énergie de la réaction qui est « libre » d'être convertie en travail. Le reste ($-\Delta H + \Delta G = -T\Delta S$ pour une évolution réversible) est échangé sous forme de chaleur. En fonctionnement réel ($i \neq 0$), l'effet Joule dans la pile et les irréversibilités aux électrodes créent de l'entropie : on récupère moins que $-\Delta G$.

\section{Tension à vide et enthalpie libre de réaction}

\begin{resultat}{Tension à vide d'une pile}
\[ \boxed{e = -\frac{\Delta_rG}{\gamma F}} \qquad\qquad \boxed{e^\circ = -\frac{\Delta_rG^\circ}{\gamma F} = E^\circ_\oplus - E^\circ_\ominus} \]
\[ e = e^\circ - \frac{RT}{\gamma F}\ln Q_r \;\underset{\SI{298}{K}}{=}\; e^\circ - \frac{0{,}06}{\gamma}\log Q_r = \frac{0{,}06}{\gamma}\log\frac{K^\circ}{Q_r}. \]
\end{resultat}

\begin{demo}
Quand l'avancement augmente de $\dd\xi$, $\gamma\,\dd\xi$ moles d'électrons traversent le circuit extérieur de $\ominus$ vers $\oplus$ : la charge $\dd q = \gamma F\,\dd\xi$ « traverse » la pile de $\ominus$ vers $\oplus$ en interne. Si $u = V_\oplus - V_\ominus$ est la tension aux bornes, la pile (en convention générateur) fournit $u\,\dd q$, donc reçoit $\delta W_{\text{élec}} = -u\,\gamma F\,\dd\xi$.

\noindent À la limite réversible ($i \to 0$), $u \to e$ et $\dd G = \delta W_{\text{élec}}$. Or, à $T$ et $P$ fixées, $\dd G = \Delta_rG\,\dd\xi$. D'où $\Delta_rG\,\dd\xi = -\gamma F e\,\dd\xi$, soit $e = -\Delta_rG/(\gamma F)$.
\end{demo}

En combinant avec la mesure $e = E_\oplus - E_\ominus$, on obtient $\Delta_rG = -\gamma F(E_\oplus - E_\ominus)$ : c'est la relation admise au chapitre~\ref{ch:redox}, maintenant justifiée.

\begin{exemple}[pile Daniell]
$e^\circ = E^\circ(\ce{Cu^{2+}/Cu}) - E^\circ(\ce{Zn^{2+}/Zn}) = 0{,}34 - (-0{,}76) = \SI{1,10}{V}$, donc $\Delta_rG^\circ = -2\times96485\times1{,}10 = \SI{-212}{kJ.mol^{-1}}$ et $\log K^\circ = \dfrac{2\times1{,}10}{0{,}059} \approx 37$.
\end{exemple}

\section{Évolution de la tension pendant la décharge}

Au cours de la décharge, $Q_r = \dfrac{[\ce{Zn^{2+}}]}{[\ce{Cu^{2+}}]}$ augmente, donc $e$ diminue. La pile est « usée » lorsque le système atteint l'équilibre :
\[ \boxed{\text{pile usée : } e = 0 \;\Longleftrightarrow\; \Delta_rG = 0 \;\Longleftrightarrow\; Q_r = K^\circ} \]

Comme $K^\circ \approx 10^{37}$, $Q_r$ reste très inférieur à $K^\circ$ pendant presque toute la décharge : $\log Q_r$ varie de quelques unités seulement, et $e$ reste quasi constante. $\Delta_rG \approx$ constante, donc $G(\xi)$ est quasi \emph{affine} : elle ne s'incurve qu'à l'extrême fin, quand \ce{Cu^{2+}} est presque épuisé.

\begin{figure}[H]
\centering
\begin{tikzpicture}
\begin{axis}[cours,width=0.47\linewidth,height=0.4\linewidth,xlabel={$q$ (C)},ylabel={$e$ (V)},
  xmin=0,xmax=1000,ymin=0.85,ymax=1.15,ytick={0.9,1.0,1.1},
  yticklabel style={/pgf/number format/fixed,/pgf/number format/precision=2},
  title={\small tension à vide},title style={yshift=-4pt}]
  \addplot[c1] table[x=q,y=E]{data/daniell.dat};
  \addplot[c3,dashed,line width=0.8pt] coordinates {(965,0.85) (965,1.15)};
  \node[c3,font=\scriptsize,anchor=east,align=right] at (axis cs:955,0.92) {$q_{\max} = \SI{965}{C}$};
\end{axis}
\end{tikzpicture}\hfill
\begin{tikzpicture}
\begin{axis}[cours,width=0.47\linewidth,height=0.4\linewidth,xlabel={$\xi$ (mmol)},ylabel={$G - G(0)$ (kJ)},
  xmin=0,xmax=5.2,ymin=-1.15,ymax=0.05,x filter/.expression={x/(2*96.485)},
  title={\small enthalpie libre},title style={yshift=-4pt}]
  \addplot[c1] table[x=q,y=G_kJ]{data/daniell.dat};
  \addplot[black!50,dashed,line width=0.8pt] coordinates {(0,0) (5,-1.061)};
  \node[font=\scriptsize,anchor=south west,align=left] at (axis cs:0.3,-0.75) {pente $\approx\Delta_rG^\circ$\\$= -2Fe^\circ$};
\end{axis}
\end{tikzpicture}
\caption{Décharge réversible d'une pile Daniell ($c_0 = \SI{0,10}{mol.L^{-1}}$ et $V = \SI{50}{mL}$ dans chaque compartiment). La tension passe de $\SI{1,10}{V}$ à $\SI{1,09}{V}$ à mi-décharge et vaut encore $\SI{1,03}{V}$ à 99\,\% ; elle ne chute qu'à l'épuisement de \ce{Cu^{2+}}.}
\label{fig:decharge}
\end{figure}

\section{Capacité d'une pile}

\begin{definition}[capacité]
La \textbf{capacité} d'une pile est la charge électrique totale qu'elle peut débiter jusqu'à épuisement du réactif limitant :
\[ \boxed{q_{\max} = \gamma\,F\,\xi_{\max}} \]
Elle s'exprime en coulombs ou, en pratique, en \textbf{ampères-heures} : $\SI{1}{A.h} = \SI{3600}{C}$.
\end{definition}

\begin{resultat}{Énergie électrique maximale}
Si $e$ reste quasi constante pendant la décharge (cas usuel, $Q_r \ll K^\circ$) :
\[ W_{\max} = -\Delta G = -\xi_{\max}\,\Delta_rG = \gamma F\,\xi_{\max}\,e = q_{\max}\,e = i\,\Delta t\,e. \]
\end{resultat}

\begin{exemple}[capacité apportée par $\SI{1}{g}$ de zinc]
$\xi_{\max} = \dfrac{m}{M} = \dfrac{1}{65{,}4} = \SI{15,3}{mmol}$ (zinc limitant), d'où
\[ q_{\max} = 2\times96485\times15{,}3\times10^{-3} = \SI{2,95e3}{C} = \frac{2950}{3600} = \SI{0,82}{A.h}. \]
L'énergie correspondante est $W_{\max} \approx q_{\max}e^\circ = 2950\times1{,}10 = \SI{3,2}{kJ}$. Sous un courant de $\SI{10}{mA}$, la pile fonctionnerait $\Delta t = q_{\max}/i \approx \SI{82}{h}$.
\end{exemple}

\begin{methode}[bilan d'une pile]
\begin{enumerate}
  \item Identifier les couples, écrire les demi-équations : réduction à la cathode ($\oplus$, plus haut potentiel), oxydation à l'anode ($\ominus$).
  \item Écrire la réaction de fonctionnement et repérer $\gamma$.
  \item Tension à vide : $e = E_\oplus - E_\ominus$ (Nernst), ou $e = -\Delta_rG/(\gamma F)$ avec des données thermodynamiques.
  \item Réactif limitant $\Rightarrow$ $\xi_{\max}$ $\Rightarrow$ capacité $q_{\max} = \gamma F\xi_{\max}$, énergie $W\approx q_{\max}e$, durée $\Delta t = q_{\max}/i$.
\end{enumerate}
\end{methode}

\begin{ordres}
\begin{itemize}
  \item Pile alcaline AA : $e \approx \SI{1,5}{V}$, capacité de l'ordre de $\SI{2,5}{A.h}$, soit environ $\SI{3,5}{W.h}$.
  \item Accumulateur lithium-ion : $\approx\SI{3,7}{V}$ par élément ; batterie de smartphone : de l'ordre de $\SI{4}{A.h}$, soit environ $\SI{15}{W.h}$.
  \item En électrolyte aqueux, la décomposition de l'eau limite en pratique la tension vers $\SI{2}{V}$ (accumulateur au plomb : $\SI{2,1}{V}$) ; les accumulateurs lithium-ion dépassent $\SI{3}{V}$ grâce à un électrolyte organique.
\end{itemize}
\end{ordres}

% ================================================================== Chapitre 9
\chapter{Questions de cours et exercices corrigés}
\label{ch:exos}

\section*{Questions de cours (réponses courtes)}
\addcontentsline{toc}{section}{Questions de cours}

\begin{enumerate}[label=\textbf{Q\arabic*.},leftmargin=2.6em,itemsep=0.5em]
  \item \emph{Définir l'état standard de référence d'un élément. Quel est-il pour l'oxygène, le brome à $\SI{298}{K}$, le carbone ?}\\
  État standard du corps simple le plus stable à $T$ : \ce{O2(g)}, \ce{Br2(l)}, C graphite.
  \item \emph{Pourquoi $\Delta_fH^\circ(\ce{O2(g)}) = 0$ mais $\Delta_fH^\circ(\ce{O3(g)})\neq0$ ?}\\
  \ce{O2(g)} est l'ESR de l'oxygène : sa réaction de formation ne transforme rien. \ce{O3} n'est pas l'ESR.
  \item \emph{Énoncer la loi de Hess.}\\
  $\Delta_rH^\circ = \sum_i\nu_i\Delta_fH^\circ_i$ ; elle découle du fait que $H$ est une fonction d'état.
  \item \emph{Donner le signe de $\Delta_rS^\circ$ pour \ce{CaCO3(s) = CaO(s) + CO2(g)}.}\\
  $\Delta\nu_{\text{gaz}} = +1$ : $\Delta_rS^\circ>0$ (on trouve $+\SI{160}{J.K^{-1}.mol^{-1}}$).
  \item \emph{Établir la relation entre $\delta S_c$ et $\Delta_rG$ à $T$ et $P$ fixées.}\\
  $\dd H = \delta Q$, $\dd S = \delta Q/T + \delta S_c$, donc $\dd G = -T\delta S_c$ ; et $\dd G = \Delta_rG\,\dd\xi$. D'où $T\delta S_c = -\Delta_rG\,\dd\xi \geq 0$.
  \item \emph{Démontrer $\Delta_rG = \Delta_rG^\circ + RT\ln Q_r$.}\\
  Injecter $\mu_i = \mu_i^\circ + RT\ln a_i$ dans $\Delta_rG = \sum\nu_i\mu_i$.
  \item \emph{Énoncer la relation de Van~'t~Hoff et la démontrer dans l'approximation d'Ellingham.}\\
  $\dd\ln K^\circ/\dd T = \Delta_rH^\circ/RT^2$ ; dériver $\ln K^\circ = -\Delta_rH^\circ/RT + \Delta_rS^\circ/R$.
  \item \emph{Effet d'un ajout de gaz inerte à $P$ fixée sur \ce{N2 + 3H2 = 2NH3} ?}\\
  $Q_r\propto n_{\text{tot}}^{2}$ augmente : évolution dans le sens indirect (défavorable).
  \item \emph{Relier $\Delta_rG^\circ$ d'une réaction rédox aux potentiels standard.}\\
  $\Delta_rG^\circ = -\gamma F(E_1^\circ - E_2^\circ)$, d'où $\log K^\circ = \gamma(E_1^\circ - E_2^\circ)/0{,}06$ à $\SI{298}{K}$.
  \item \emph{Établir l'inégalité $\Delta G\leq W_{\text{élec}}$.}\\
  $\dd H = \delta Q + \delta W_{\text{élec}}$ (monobare), $\dd S = \delta Q/T + \delta S_c$, donc $\dd G = \delta W_{\text{élec}} - T\delta S_c$.
  \item \emph{Que valent $e$, $\Delta_rG$ et $Q_r$ pour une pile usée ?}\\
  $e = 0$, $\Delta_rG = 0$, $Q_r = K^\circ$ : la pile est à l'équilibre.
  \item \emph{Définir la capacité d'une pile. Convertir $\SI{1}{A.h}$ en coulombs.}\\
  $q_{\max} = \gamma F\xi_{\max}$ ; $\SI{1}{A.h} = \SI{3600}{C}$.
\end{enumerate}

\section*{Données pour les exercices}
\addcontentsline{toc}{section}{Données pour les exercices}
Les données thermodynamiques à $\SI{298}{K}$ sont rassemblées dans le tableau~\ref{tab:donnees} (chapitre~\ref{ch:formulaire}). On utilise l'approximation d'Ellingham, $R = \SI{8,314}{J.K^{-1}.mol^{-1}}$ et $F = \SI{96485}{C.mol^{-1}}$.

% ------------------------------------------------------------------
\exercice{Synthèse de l'ammoniac}{2}
\label{exo:nh3}
On étudie \ce{N2(g) + 3 H2(g) = 2 NH3(g)}.
\begin{enumerate}
  \item Prévoir le signe de $\Delta_rS^\circ$, puis calculer $\Delta_rH^\circ$ et $\Delta_rS^\circ$.
  \item Exprimer $\Delta_rG^\circ(T)$. Calculer la température d'inversion et $K^\circ(\SI{298}{K})$.
  \item Calculer $K^\circ$ à $\SI{723}{K}$ de deux façons.
  \item Dans un réacteur à $\SI{723}{K}$, les pressions partielles sont $P_{\ce{N2}} = \SI{50}{bar}$, $P_{\ce{H2}} = \SI{150}{bar}$, $P_{\ce{NH3}} = \SI{20}{bar}$. Dans quel sens le système évolue-t-il ?
  \item Justifier les conditions industrielles : $\SI{450}{\celsius}$, $\SI{200}{bar}$, catalyseur au fer, liquéfaction de \ce{NH3}.
\end{enumerate}

\correction
\begin{enumerate}
  \item $\Delta\nu_{\text{gaz}} = -2$ : $\Delta_rS^\circ<0$. $\Delta_rH^\circ = 2\Delta_fH^\circ(\ce{NH3}) = \SI{-91,8}{kJ.mol^{-1}}$ (corps simples dans leur ESR). $\Delta_rS^\circ = 2(192{,}8) - 191{,}6 - 3(130{,}7) = \SI{-198,1}{J.K^{-1}.mol^{-1}}$.
  \item $\Delta_rG^\circ(T) = -91{,}8\times10^3 + 198{,}1\,T$ (en $\mathrm{J\cdot mol^{-1}}$). $T_i = 91{,}8\times10^3/198{,}1 = \SI{463}{K}$. À $\SI{298}{K}$ : $\Delta_rG^\circ = \SI{-32,7}{kJ.mol^{-1}}$, $K^\circ = \exp(32{,}7\times10^3/(8{,}314\times298)) = \num{5,4e5}$.
  \item Directement : $\Delta_rG^\circ(723) = \SI{51,4}{kJ.mol^{-1}}$, $K^\circ = \num{1,9e-4}$. Par Van~'t~Hoff intégrée depuis $\SI{298}{K}$ : même valeur (voir chapitre~\ref{ch:vanthoff}).
  \item $Q_r = \dfrac{(P_{\ce{NH3}}/P^\circ)^2}{(P_{\ce{N2}}/P^\circ)(P_{\ce{H2}}/P^\circ)^3} = \dfrac{20^2}{50\times150^3} = \num{2,4e-6} < K^\circ$. Donc $\Delta_rG = RT\ln(Q_r/K^\circ) = \SI{-26}{kJ.mol^{-1}} < 0$ : évolution dans le \textbf{sens direct}, formation d'ammoniac.
  \item Réaction exothermique : $K^\circ$ diminue avec $T$, mais en dessous de $\SI{400}{\celsius}$ la réaction est trop lente même catalysée $\Rightarrow$ compromis vers $\SI{450}{\celsius}$. $\Delta\nu_{\text{gaz}}<0$ : une haute pression diminue $Q_r$ et favorise la synthèse. Le catalyseur accélère sans changer $K^\circ$. La liquéfaction retire \ce{NH3}, ce qui diminue $Q_r$ et permet de recycler \ce{N2} et \ce{H2}.
\end{enumerate}

% ------------------------------------------------------------------
\exercice{Effets thermiques : flamme et calorimétrie}{2}
\label{exo:gaz}
\textbf{A. Flamme dans le dioxygène pur.} On reprend la combustion \ce{CH4(g) + 2 O2(g) -> CO2(g) + 2 H2O(g)} ($\Delta_rH^\circ = \SI{-802,5}{kJ.mol^{-1}}$), mais avec un mélange stœchiométrique dans \ce{O2} pur, initialement à $\SI{298}{K}$.
\begin{enumerate}
  \item Calculer la température de flamme adiabatique et la comparer à celle obtenue dans l'air ($\approx\SI{2680}{K}$).
  \item La température réelle d'une flamme oxy-méthane est de l'ordre de $\SI{3000}{K}$. Commenter.
\end{enumerate}
\textbf{B. Mesure d'une enthalpie de réaction.} Dans un calorimètre de capacité thermique $C_{\text{cal}} = \SI{120}{J.K^{-1}}$, on mélange $\SI{100}{mL}$ d'acide chlorhydrique à $\SI{1,00}{mol.L^{-1}}$ et $\SI{100}{mL}$ de soude à $\SI{1,00}{mol.L^{-1}}$, tous deux à $\SI{20,0}{\celsius}$. On mesure une température finale de $\SI{25,8}{\celsius}$. On assimile les solutions à de l'eau ($c_{\text{eau}} = \SI{4,18}{J.g^{-1}.K^{-1}}$, $\rho = \SI{1}{g.mL^{-1}}$).
\begin{enumerate}[resume]
  \item Écrire la réaction et justifier qu'elle est totale.
  \item En déduire $\Delta_rH^\circ$. Comparer à la valeur tabulée calculée avec $\Delta_fH^\circ(\ce{H2O(l)}) = \SI{-285,8}{kJ.mol^{-1}}$, $\Delta_fH^\circ(\ce{HO-}) = \SI{-230,0}{kJ.mol^{-1}}$ et la convention $\Delta_fH^\circ(\ce{H+}) = 0$.
\end{enumerate}

\correction
\begin{enumerate}
  \item État final : \ce{CO2 + 2 H2O(g)}, $\sum n_jC_{p,m,j} = 37{,}1 + 2\times33{,}6 = \SI{104,3}{J.K^{-1}}$. $T_F = 298 + \dfrac{802{,}5\times10^3}{104{,}3} \approx \SI{8000}{K}$. Sans le diazote, il y a trois fois moins de matière à chauffer : la température calculée est trois fois plus élevée.
  \item La valeur calculée est absurde : à ces températures, \ce{CO2} et \ce{H2O} se dissocient (réactions très endothermiques qui « absorbent » l'énergie) et les $C_p$ augmentent. L'hypothèse « combustion totale, $C_p$ constants » n'est plus valable. Elle reste correcte qualitativement : la flamme dans \ce{O2} pur est bien plus chaude (d'où son usage en soudage et découpage).
  \item \ce{H+ + HO- -> H2O}, de constante $K^\circ = 1/K_e = 10^{14}$ : réaction quasi totale. $\xi_F = \SI{0,100}{mol}$.
  \item Calorimètre adiabatique, monobare : $\Delta H = 0 = \xi_F\Delta_rH^\circ + (m_{\text{eau}}c_{\text{eau}} + C_{\text{cal}})(T_F - T_I)$, avec $m_{\text{eau}}c_{\text{eau}} = 200\times4{,}18 = \SI{836}{J.K^{-1}}$ :
  \[ \Delta_rH^\circ = -\frac{(836 + 120)\times5{,}8}{0{,}100} = \SI{-55,4}{kJ.mol^{-1}}. \]
  Valeur tabulée : $\Delta_rH^\circ = -285{,}8 - (-230{,}0) - 0 = \SI{-55,8}{kJ.mol^{-1}}$. L'écart (moins de 1\,\%) vient des pertes thermiques et de l'assimilation des solutions à de l'eau pure.
\end{enumerate}

% ------------------------------------------------------------------
\exercice{Équilibre de Boudouard}{2}
\label{exo:boudouard}
On étudie \ce{C(graphite) + CO2(g) = 2 CO(g)}, réaction clé du haut fourneau.
\begin{enumerate}
  \item Calculer $\Delta_rH^\circ$ et $\Delta_rS^\circ$ ; commenter leurs signes. Calculer la température d'inversion.
  \item Calculer $K^\circ$ à $\SI{1000}{K}$ et à $\SI{800}{K}$.
  \item On place $\SI{1}{mol}$ de \ce{CO2} au contact d'un excès de graphite à $\SI{1000}{K}$, sous $P = \SI{1}{bar}$. Déterminer la composition à l'équilibre.
  \item Même question sous $P = \SI{10}{bar}$, puis sous $\SI{1}{bar}$ en présence de $\SI{4}{mol}$ de \ce{N2} (le gaz du haut fourneau provient de l'air). Interpréter.
\end{enumerate}

\correction
\begin{enumerate}
  \item $\Delta_rH^\circ = 2(-110{,}5) - (-393{,}5) = \SI{+172,5}{kJ.mol^{-1}}$ : endothermique. $\Delta_rS^\circ = 2(197{,}7) - 213{,}8 - 5{,}7 = \SI{+175,9}{J.K^{-1}.mol^{-1}}$ : positive, cohérente avec $\Delta\nu_{\text{gaz}} = +1$. $T_i = 172{,}5\times10^3/175{,}9 = \SI{981}{K}$.
  \item $K^\circ(1000) = \exp\left(-\dfrac{172{,}5\times10^3 - 1000\times175{,}9}{8{,}314\times1000}\right) = 1{,}50$ ; $K^\circ(800) = \num{8,4e-3}$. Réaction endothermique : $K^\circ$ croît avec $T$ (Van~'t~Hoff).
  \item Le graphite est en excès (activité 1). Avec $n_{\ce{CO2}} = 1-\xi$, $n_{\ce{CO}} = 2\xi$, $n_{\text{tot,gaz}} = 1+\xi$ :
  \[ K^\circ = \frac{(P_{\ce{CO}}/P^\circ)^2}{P_{\ce{CO2}}/P^\circ} = \frac{4\xi^2}{(1-\xi)(1+\xi)}\,\frac{P}{P^\circ} \;\Longrightarrow\; \xi_{\text{éq}} = \sqrt{\frac{K^\circ}{K^\circ + 4P/P^\circ}} = \sqrt{\frac{1{,}50}{5{,}50}} = \SI{0,52}{mol}. \]
  Composition : $\SI{0,48}{mol}$ de \ce{CO2}, $\SI{1,05}{mol}$ de \ce{CO}, soit $x_{\ce{CO}} = 0{,}69$.
  \item Sous $\SI{10}{bar}$ : $\xi_{\text{éq}} = \sqrt{1{,}50/41{,}5} = \SI{0,19}{mol}$ ($x_{\ce{CO}} = 0{,}32$) : $\Delta\nu_{\text{gaz}} > 0$, une hausse de pression défavorise la formation de CO. Avec $\SI{4}{mol}$ de \ce{N2} à $\SI{1}{bar}$ : $\dfrac{4\xi^2}{(1-\xi)(5+\xi)} = 1{,}50$ donne (résolution numérique) $\xi_{\text{éq}} = \SI{0,74}{mol}$ : la dilution par un inerte à $P$ fixée équivaut à une baisse de pression et favorise ici le sens direct. Le haut fourneau produit donc efficacement le réducteur CO à haute température et sous pression modérée.
\end{enumerate}

% ------------------------------------------------------------------
\exercice{Pile à combustible à hydrogène}{2}
\label{exo:pac}
Dans une pile à combustible à électrolyte acide, le dihydrogène est oxydé à une électrode et le dioxygène réduit à l'autre, avec formation d'eau liquide.
\begin{enumerate}
  \item Écrire les demi-équations, préciser l'anode et la cathode, et écrire la réaction de fonctionnement pour une mole de \ce{H2}.
  \item Calculer la tension à vide standard $e^\circ$ à $\SI{298}{K}$.
  \item Quelle énergie électrique maximale peut-on obtenir avec $\SI{1}{kg}$ de \ce{H2} ? Quelle capacité (en A$\cdot$h) cela représente-t-il ?
  \item On définit le rendement maximal $\eta = \Delta_rG^\circ/\Delta_rH^\circ$. Le calculer et l'interpréter.
  \item Calculer $\dd e^\circ/\dd T$, puis $e^\circ$ à $\SI{80}{\celsius}$ (température de fonctionnement usuelle).
  \item La cathode est alimentée par de l'air ($P_{\ce{O2}} = \SI{0,2}{bar}$) et l'anode par \ce{H2} sous $\SI{1}{bar}$. Calculer $e$ à $\SI{298}{K}$.
\end{enumerate}

\correction
\begin{enumerate}
  \item Anode ($\ominus$, oxydation) : \ce{H2 -> 2H+ + 2e^-}. Cathode ($\oplus$, réduction) : \ce{O2 + 4H+ + 4e^- -> 2H2O}. Bilan pour une mole de \ce{H2} : \ce{H2(g) + 1/2 O2(g) -> H2O(l)}, $\gamma = 2$.
  \item $\Delta_rG^\circ = \Delta_fG^\circ(\ce{H2O(l)}) = \SI{-237,1}{kJ.mol^{-1}}$, donc $e^\circ = \dfrac{237{,}1\times10^3}{2\times96485} = \SI{1,23}{V}$ (on retrouve $E^\circ(\ce{O2/H2O}) - E^\circ(\ce{H+/H2})$).
  \item $n = 1000/2{,}016 = \SI{496}{mol}$ ; $W_{\max} = n\times237{,}1 = \SI{118}{MJ} \approx \SI{33}{kWh}$. Capacité : $q_{\max} = 2Fn = \SI{9,6e7}{C} \approx \SI{2,7e4}{A.h}$.
  \item $\eta = 237{,}1/285{,}8 = 0{,}83$. Au mieux, 83\,\% de l'énergie de la réaction devient électrique ; le reste ($-T\Delta_rS^\circ$, avec $\Delta_rS^\circ<0$) est forcément cédé sous forme de chaleur. Ce rendement n'est pas limité par Carnot, contrairement à celui d'un moteur thermique.
  \item $e^\circ = -\dfrac{\Delta_rH^\circ - T\Delta_rS^\circ}{2F}$ donc $\dfrac{\dd e^\circ}{\dd T} = \dfrac{\Delta_rS^\circ}{2F}$, avec $\Delta_rS^\circ = 70{,}0 - 130{,}7 - \tfrac12(205{,}2) = \SI{-163,3}{J.K^{-1}.mol^{-1}}$ : $\dd e^\circ/\dd T = \SI{-0,85}{mV.K^{-1}}$. À $\SI{353}{K}$ : $e^\circ = 1{,}229 - 55\times0{,}85\times10^{-3} = \SI{1,18}{V}$.
  \item $Q_r = \dfrac{1}{(P_{\ce{H2}}/P^\circ)(P_{\ce{O2}}/P^\circ)^{1/2}}$, donc $e = e^\circ - \dfrac{0{,}059}{2}\log Q_r = 1{,}229 + \dfrac{0{,}059}{4}\log 0{,}2 = \SI{1,22}{V}$. L'effet de la dilution du dioxygène est faible.
\end{enumerate}

% ------------------------------------------------------------------
\exercice{Degrés d'oxydation du cuivre}{2}
\label{exo:cuivre}
Données : $E^\circ(\ce{Cu^{2+}/Cu}) = \SI{0,34}{V}$, $E^\circ(\ce{Cu^{2+}/Cu+}) = \SI{0,16}{V}$.
\begin{enumerate}
  \item Calculer $E^\circ(\ce{Cu+/Cu})$.
  \item Écrire la réaction de dismutation de \ce{Cu+} et calculer $\Delta_rG^\circ$ et $K^\circ$.
  \item On prépare une solution de \ce{Cu+} à $\SI{0,010}{mol.L^{-1}}$. Que se passe-t-il ? Calculer la concentration résiduelle en \ce{Cu+}.
\end{enumerate}

\correction
\begin{enumerate}
  \item \ce{Cu^{2+} + 2e^- = Cu} est la somme de \ce{Cu^{2+} + e^- = Cu+} et \ce{Cu+ + e^- = Cu} : $-2F(0{,}34) = -F(0{,}16) - F\,E^\circ(\ce{Cu+/Cu})$, donc $E^\circ(\ce{Cu+/Cu}) = 2\times0{,}34 - 0{,}16 = \SI{0,52}{V}$.
  \item \ce{2 Cu+ = Cu^{2+} + Cu(s)} ($\gamma = 1$) : \ce{Cu+} est oxydant du couple \ce{Cu+/Cu} ($\SI{0,52}{V}$) et réducteur du couple \ce{Cu^{2+}/Cu+} ($\SI{0,16}{V}$). $\Delta_rG^\circ = -F(0{,}52 - 0{,}16) = \SI{-34,7}{kJ.mol^{-1}}$, $K^\circ = 10^{0{,}36/0{,}06} = 10^6$.
  \item $K^\circ$ est grand : dismutation quasi totale, formation de cuivre métallique et de \ce{Cu^{2+}} à environ $\SI{0,005}{mol.L^{-1}}$. Alors $\dfrac{[\ce{Cu^{2+}}]c^\circ}{[\ce{Cu+}]^2} = 10^6$ donne $[\ce{Cu+}] = \sqrt{0{,}005\times10^{-6}} = \SI{7e-5}{mol.L^{-1}}$. \ce{Cu+} n'est pas stable en solution aqueuse : la règle du gamma le montre aussi, car l'oxydant \ce{Cu+} est au-dessus du réducteur \ce{Cu+} sur l'échelle des potentiels.
\end{enumerate}

% ------------------------------------------------------------------
\exercice{Pile bouton argent--zinc}{2}
\label{exo:bouton}
Une pile bouton de montre fonctionne selon \ce{Zn(s) + Ag2O(s) -> ZnO(s) + 2 Ag(s)} dans un électrolyte de potasse concentrée.

\noindent Données : $\Delta_fG^\circ(\ce{ZnO(s)}) = \SI{-320,5}{kJ.mol^{-1}}$ ; $\Delta_fG^\circ(\ce{Ag2O(s)}) = \SI{-11,2}{kJ.mol^{-1}}$ ; $M(\ce{Ag2O}) = \SI{231,7}{g.mol^{-1}}$ ; $M(\ce{Zn}) = \SI{65,4}{g.mol^{-1}}$.
\begin{enumerate}
  \item Écrire les demi-équations en milieu basique ; identifier l'anode et la cathode.
  \item Calculer la tension à vide standard. Pourquoi la tension à vide reste-t-elle constante pendant toute la décharge ?
  \item La pile contient $\SI{0,50}{g}$ d'oxyde d'argent (réactif limitant). Calculer sa capacité en mA$\cdot$h et la masse minimale de zinc nécessaire.
  \item Une montre consomme $\SI{2}{\micro A}$ en moyenne. Estimer la durée de vie de la pile.
\end{enumerate}

\correction
\begin{enumerate}
  \item Anode ($\ominus$) : \ce{Zn + 2HO- -> ZnO + H2O + 2e^-}. Cathode ($\oplus$) : \ce{Ag2O + H2O + 2e^- -> 2Ag + 2HO-}. Bilan avec $\gamma = 2$.
  \item $\Delta_rG^\circ = -320{,}5 - (-11{,}2) = \SI{-309,3}{kJ.mol^{-1}}$, $e^\circ = \dfrac{309{,}3\times10^3}{2\times96485} = \SI{1,60}{V}$. Tous les constituants du bilan sont des solides purs ($a = 1$) : $Q_r = 1$ en permanence, donc $e = e^\circ$ \emph{tout au long} de la décharge (les ions \ce{HO-} et l'eau s'éliminent du bilan). C'est une qualité recherchée pour l'horlogerie.
  \item $n(\ce{Ag2O}) = 0{,}50/231{,}7 = \SI{2,16}{mmol} = \xi_{\max}$ ; $q_{\max} = 2F\xi_{\max} = \SI{416}{C} = \SI{116}{mA.h}$. Zinc : $\SI{2,16}{mmol}\times65{,}4 = \SI{141}{mg}$ au minimum. Énergie : $q_{\max}e^\circ = \SI{667}{J}$.
  \item $\Delta t = q_{\max}/i = 416/(2\times10^{-6}) = \SI{2,1e8}{s} \approx 6{,}6$ ans (en pratique un peu moins, à cause de l'autodécharge).
\end{enumerate}

% ================================================================== Chapitre 10
\chapter{Formulaire et données}
\label{ch:formulaire}

\section*{Formulaire}
\addcontentsline{toc}{section}{Formulaire}

\renewcommand{\arraystretch}{1.55}
\begin{xltabular}{\linewidth}{>{\raggedright\arraybackslash}p{4.3cm} X}
\toprule
\multicolumn{2}{l}{\sffamily\bfseries\color{navy} Potentiel chimique et enthalpie libre} \\ \midrule
Enthalpie libre & $G = H - TS = U + PV - TS$ \\
Identité (composition fixée) & $\dd G = V\,\dd P - S\,\dd T$ \\
Mélange & $\dd G = V\,\dd P - S\,\dd T + \sum_i\mu_i\,\dd n_i$, \quad $G = \sum_in_i\mu_i$ \\
Potentiel chimique & $\mu_i = \left(\partial G/\partial n_i\right)_{T,P,n_{j\neq i}} = \mu_i^\circ(T) + RT\ln a_i$ \\
Activités & gaz : $P_i/P^\circ$ ; soluté : $c_i/c^\circ$ ; solvant, solide ou liquide pur : $1$ \\ \midrule
\multicolumn{2}{l}{\sffamily\bfseries\color{navy} Grandeurs de réaction} \\ \midrule
Grandeur de réaction & $\Delta_rX = (\partial X/\partial\xi)_{T,P} = \sum_i\nu_iX_{m,i}$ \quad ; \quad $\Delta_rG = \sum_i\nu_i\mu_i$ \\
Grandeur standard & $\Delta_rX^\circ(T) = \sum_i\nu_iX^\circ_{m,i}$ \quad ; \quad $\Delta_rG^\circ = \Delta_rH^\circ - T\Delta_rS^\circ$ \\
Loi de Hess & $\Delta_rH^\circ = \sum_i\nu_i\Delta_fH^\circ_i$, \quad $\Delta_rG^\circ = \sum_i\nu_i\Delta_fG^\circ_i$, \quad $\Delta_rS^\circ = \sum_i\nu_iS^\circ_{m,i}$ \\
ESR & $\Delta_fX^\circ(\text{corps simple dans son ESR}) = 0$ \quad (mais $S^\circ_m \neq 0$) \\
Combinaison linéaire & $(3) = a(1) + b(2) \Rightarrow \Delta_rX^\circ_3 = a\Delta_rX^\circ_1 + b\Delta_rX^\circ_2$, \quad $K^\circ_3 = (K^\circ_1)^a(K^\circ_2)^b$ \\
Ellingham & $\Delta_rH^\circ$, $\Delta_rS^\circ$ constantes hors changement d'état ; $T_i = \Delta_rH^\circ/\Delta_rS^\circ$ \\ \midrule
\multicolumn{2}{l}{\sffamily\bfseries\color{navy} Premier principe} \\ \midrule
Monobare isotherme & $Q = \Delta H = \xi_F\,\Delta_rH^\circ$ \\
Monobare adiabatique & $\xi_F\,\Delta_rH^\circ(T_I) + \big(\sum_jn_jC^\circ_{p,m,j}\big)(T_F - T_I) = 0$ \\ \midrule
\multicolumn{2}{l}{\sffamily\bfseries\color{navy} Évolution et équilibre ($T$, $P$ fixées)} \\ \midrule
Création d'entropie & $T\,\delta S_c = -\Delta_rG\,\dd\xi \geq 0$ \quad ; \quad $\dd G = -T\,\delta S_c \leq 0$ \\
Enthalpie libre de réaction & $\Delta_rG = \Delta_rG^\circ + RT\ln Q_r = RT\ln(Q_r/K^\circ)$, \quad $Q_r = \prod_ia_i^{\nu_i}$ \\
Constante d'équilibre & $K^\circ = \exp(-\Delta_rG^\circ/RT)$ \quad ; \quad équilibre : $Q_r = K^\circ$, $\Delta_rG = 0$ \\
Sens d'évolution & $\Delta_rG<0 \Leftrightarrow Q_r<K^\circ$ : sens direct \\
Van~'t~Hoff & $\dfrac{\dd\ln K^\circ}{\dd T} = \dfrac{\Delta_rH^\circ}{RT^2}$ \quad ; \quad $\ln\dfrac{K^\circ_2}{K^\circ_1} = -\dfrac{\Delta_rH^\circ}{R}\left(\dfrac1{T_2} - \dfrac1{T_1}\right)$ \\
Gaz : effet de $P$ & $Q_r = \prod n_i^{\nu_i}\big(P/(n_{\text{tot,gaz}}P^\circ)\big)^{\Delta\nu_{\text{gaz}}}$ \\ \midrule
\multicolumn{2}{l}{\sffamily\bfseries\color{navy} Oxydoréduction et piles} \\ \midrule
Nernst & $E = E^\circ + \dfrac{RT}{\gamma F}\ln\dfrac{a_{\mathrm{Ox}}^\alpha}{a_{\mathrm{Red}}^\beta} \approx E^\circ + \dfrac{0{,}06}{\gamma}\log\dfrac{a_{\mathrm{Ox}}^\alpha}{a_{\mathrm{Red}}^\beta}$ \\
Réaction rédox & $\Delta_rG = -\gamma F(E_1 - E_2)$, \quad $\Delta_rG^\circ = -\gamma F(E^\circ_1 - E^\circ_2)$ \\
Constante & $\log K^\circ = \gamma(E^\circ_1 - E^\circ_2)/0{,}06$ \quad à $\SI{298}{K}$ \\
Demi-équation & $\Delta_{1/2}G^\circ = -\gamma FE^\circ$ \quad (on additionne les $\Delta G^\circ$, jamais les $E^\circ$) \\
Travail électrique & $\Delta G \leq W_{\text{élec}}$ \quad ; \quad $W_{\max,\text{fourni}} = -\Delta G$ \\
Tension à vide & $e = E_\oplus - E_\ominus = -\Delta_rG/(\gamma F) = e^\circ - \dfrac{0{,}06}{\gamma}\log Q_r$ \quad ; \quad pile usée : $e = 0$ \\
Capacité & $q_{\max} = \gamma F\xi_{\max}$, \quad $\SI{1}{A.h} = \SI{3600}{C}$, \quad $W\approx q_{\max}\,e$ \\
\bottomrule
\end{xltabular}

\begin{pieges}
\begin{enumerate}[leftmargin=1.8em]
  \item Oublier les entropies molaires des corps simples dans $\Delta_rS^\circ$.
  \item Mélanger kJ et J dans $\Delta_rG^\circ = \Delta_rH^\circ - T\Delta_rS^\circ$.
  \item Utiliser le signe de $\Delta_rG^\circ$ (au lieu de $\Delta_rG$) pour prévoir le sens d'évolution.
  \item Oublier les espèces spectatrices (\ce{N2} de l'air, eau, calorimètre) dans le calcul d'une température finale.
  \item Additionner des potentiels standard au lieu des enthalpies libres $-\gamma FE^\circ$.
  \item Écrire $\Delta_rG^\circ = +\gamma FE^\circ$ : le signe moins est indispensable (une pile, spontanée, a $e^\circ > 0$ et $\Delta_rG^\circ<0$).
  \item Écrire $G = U - PV - TS$ : c'est $G = U + PV - TS$.
  \item Croire qu'ajouter un gaz inerte n'a jamais d'effet : c'est vrai à $V$ fixé, pas à $P$ fixée.
\end{enumerate}
\end{pieges}

\section*{Données thermodynamiques à 298 K}
\addcontentsline{toc}{section}{Données thermodynamiques}

\begin{table}[H]
\centering
\renewcommand{\arraystretch}{1.18}
\small
\begin{tabular}{lrrrr}
\toprule
Constituant & $\Delta_fH^\circ$ ($\mathrm{kJ\cdot mol^{-1}}$) & $S^\circ_m$ ($\mathrm{J\cdot K^{-1}\cdot mol^{-1}}$) & $\Delta_fG^\circ$ ($\mathrm{kJ\cdot mol^{-1}}$) & $C^\circ_{p,m}$ ($\mathrm{J\cdot K^{-1}\cdot mol^{-1}}$) \\ \midrule
\ce{N2(g)}      & 0       & 191,6 & 0       & 29,1 \\
\ce{H2(g)}      & 0       & 130,7 & 0       & 28,8 \\
\ce{O2(g)}      & 0       & 205,2 & 0       & 29,4 \\
\ce{Cl2(g)}     & 0       & 223,1 & 0       & 33,9 \\
C (graphite)    & 0       & 5,7   & 0       & 8,5 \\
C (diamant)     & 1,9     & 2,4   & 2,9     & 6,1 \\
\ce{NH3(g)}     & $-45{,}9$   & 192,8 & $-16{,}4$  & 35,1 \\
\ce{CH4(g)}     & $-74{,}6$   & 186,3 & $-50{,}5$  & 35,7 \\
\ce{CH3Cl(g)}   & $-81{,}9$   & 234,6 & --      & -- \\
\ce{HCl(g)}     & $-92{,}3$   & 186,9 & $-95{,}3$  & 29,1 \\
\ce{CO(g)}      & $-110{,}5$  & 197,7 & $-137{,}2$ & 29,1 \\
\ce{CO2(g)}     & $-393{,}5$  & 213,8 & $-394{,}4$ & 37,1 \\
\ce{H2O(l)}     & $-285{,}8$  & 70,0  & $-237{,}1$ & 75,3 \\
\ce{H2O(g)}     & $-241{,}8$  & 188,8 & $-228{,}6$ & 33,6 \\
\ce{CaCO3(s)}   & $-1207{,}6$ & 91,7  & $-1129{,}1$ & 83,5 \\
\ce{CaO(s)}     & $-634{,}9$  & 38,1  & $-603{,}3$ & 42,0 \\
\ce{ZnO(s)}     & $-350{,}5$  & 43,7  & $-320{,}5$ & -- \\
\ce{Ag2O(s)}    & $-31{,}1$   & 121,3 & $-11{,}2$  & -- \\
\ce{HO-(aq)}    & $-230{,}0$  & --    & --      & -- \\
\bottomrule
\end{tabular}
\caption{Données thermodynamiques à $\SI{298}{K}$ (valeurs usuelles des tables). Convention pour les ions : $\Delta_fH^\circ(\ce{H+(aq)}) = 0$.}
\label{tab:donnees}
\end{table}

\begin{table}[H]
\centering
\renewcommand{\arraystretch}{1.18}
\small
\begin{tabular}{lr@{\hspace{2.5em}}lr}
\toprule
Couple & $E^\circ$ (V) & Couple & $E^\circ$ (V) \\ \midrule
\ce{O2/H2O}       & 1,23  & \ce{Cu^{2+}/Cu+}  & 0,16 \\
\ce{Hg2^{2+}/Hg}  & 0,797 & \ce{H+/H2}        & 0 (référence) \\
\ce{Fe^{3+}/Fe^{2+}} & 0,77 & \ce{Fe^{2+}/Fe}  & $-0{,}44$ \\
\ce{Cu^{2+}/Cu}   & 0,34  & \ce{Zn^{2+}/Zn}   & $-0{,}76$ \\
\bottomrule
\end{tabular}
\caption{Potentiels standard à $\SI{298}{K}$ utilisés dans ce cours.}
\label{tab:potentiels}
\end{table}


\end{document}
